\chapter{Teoria}
\label{cap:teoria}

\section{Variables i constants}

Quan escrivim a la pissarra una expressió com $x+y=y+x$, a molta gent se
li acut de seguida la propietat commutativa de la suma. És clar que sí, i a més, d'aquesta manera no hem d'escriure a la pissarra que
l'ordre dels sumands no altera el resultat de la suma. Hem aconseguit
sintetitzar una idea amb una expressió simbòlica fent ús de dues
variables $x$ i $y$. Si després escrivim $x+2=2+x$, tothom s'adona que
aquesta igualtat és un cas particular de l'anterior substituint $y$ per
la constant el nombre $2$. Podríem afegir molts més exemples, però no cal, tots recordeu el plantejament de problemes per equacions i la
seva resolució. Tot això seria molt més difícil de fer
sense l'ajut de variables i constants. En resum, heu de tenir molt present
que en el llenguatge de les matemàtiques és imprescindible l'ús
de variables i constants.

Com hem dit, la formalització d'enunciats amb l'ajut de variables i
constants es fa indispensable per a continuar endavant. En l'exemple següent es veu com es fa molt més senzill i alhora més rigorós
treballar amb enunciats matemàtics si introduïm variables i
constants. És habitual prendre com a variables les últimes lletres
de l'abecedari i com a constants les primeres, però això és
irrellevant.

\begin{exemple}
Amb l'ajut de variables i constants expressa els següents enunciats:

\begin{enumerate}
\item La mitja geomètrica de dos nombres positius és menor o igual
que la mitja aritmètica.

Si $x,y$ són dos nombres positius, aleshores es compleix $\sqrt{x\cdot y}\leq\dfrac{x+y}{2}$.

\item El quadrat d'un nombre senar és senar.

Si $x$ és un nombre senar, aleshores $x^{2}$ és senar.

\item Donat un nombre real positiu, busqueu els nombres reals el quadrat
dels quals és més petit o igual que el nombre donat.

Sigui $a>0$, aleshores hem de resoldre la inequació $x^{2}\leq a$.

\item Trobar tres nombres parells consecutius que sumen $12$.

Busqueu un nombre parell $x$ tal que $x+x+2+x+4=12$.
\end{enumerate}
\end{exemple}

\section{Proposicions i predicats{}}

És evident que l'enunciat \enquote{si els elefants volen, existeix el nombre $\pi $} no és matemàtic perquè hi apareixen termes que no estan
definits dins de les matemàtiques; en canvi, \enquote{$161$ és múltiple
de $7$} sí que és un enunciat matemàtic. Aquí no volem
definir en rigor què és un enunciat matemàtic, només ens
interessa que tingueu una idea intuïtiva clara sobre com distingir un
enunciat matemàtic d'aquell que no ho és, i això creiem que ho
podem assumir sense cap problema.

Hi han enunciats de les matemàtiques que són proposicions i altres,
predicats. Una \textbf{proposició} és un enunciat que és o
vertader o fals. Per exemple, els enunciats següents són
proposicions vertaderes:

\begin{itemize}
\item Si un cercle té un radi de $r$ unitats, aleshores la seva àrea
és $\pi r^{2}$ unitats quadrades.

\item Tot nombre parell és divisible per $2$.

\item $2\in\mathbb{Z}$, que es llegeix \enquote{$2$ pertany als nombres enters}.

\item $\sqrt{2}\notin\mathbb{Q}$, que es llegeix \enquote{$\sqrt{2}$ no pertany als
nombres racionals}.

\item $\mathbb{N}\subseteq\mathbb{Z}$, que es llegeix \enquote{els nombres naturals
formen un subconjunt dels nombres enters}.

\item El conjunt $\left\{ a,b,c\right\} $ té 3 elements.

\item Alguns triangles rectangles són isòsceles.
\end{itemize}

En canvi, les proposicions següents són falses:

\begin{itemize}
\item $1=2$.

\item $\sqrt{2}\notin\mathbb{R}$.

\item Per a tot nombre real $x$, es té que $x^{2}+1\leq0$.

\item $\left\{ 1,2,3\right\} \cap\mathbb{N}=\varnothing$, que es llegeix \enquote{la
intersecció dels conjunts $\left\{ 1,2,3\right\} $ i $\mathbb{N}$ és
el conjunt buit}.

\item La suma dels angles de qualsevol quadrilàter val 180$^\circ$
.

\item Si $x=2$, aleshores $x^{2}=9$.

\item 6 és un nombre primer.
\end{itemize}

Volem observar que hi ha proposicions que tenen variables. Per exemple,
l'enunciat \enquote{Si $x$ és un nombre enter, $2x+1$ és senar} és una
proposició vertadera. L'enunciat \enquote{Per a tots els nombres reals $x,y$ es
compleix $(x-y)(x+y)=x^{2}-y^{2}$} és una proposició vertadera. També hi ha proposicions falses amb variables com, per exemple, \enquote{Per a tot
nombre real $x$ es té $\sqrt{x^{2}}=x$}.

També hi ha enunciats matemàtics que encara no sabem demostrar ni
refutar. Aquests enunciats s'anomenen \textbf{conjectures}. Des del punt de
vista de la lògica clàssica continuen sent proposicions, perquè tenen
un valor de veritat determinat, encara que nosaltres no el coneguem. Per
exemple, la conjectura de Goldbach diu: tot enter parell superior a 2 és una
suma de dos nombres primers. També tenim l'antic teorema de Fermat, segons el
qual per a tots els nombres naturals $a,b,c,n$, amb $n>2$, es té
$a^{n}+b^{n}\neq c^{n}$. Durant segles va ser una conjectura; el 1994, Andrew
Wiles en va donar una demostració.

Finalment, els enunciats següents no són proposicions:

\begin{itemize}
\item $x$ és un nombre parell.

\item Una recta paral·lela a la recta d'equació $x+y=1.$

\item Els nombres enters parells més grans que un nombre irracional donat.

\item La funció $f$ és la funció inversa de la funció $g$.

\item $2x+1$.
\end{itemize}

Els \textbf{predicats} són enunciats oberts perquè contenen
variables i són vertaders o falsos depenent dels valors assignats a les
variables. Aquí denotem els predicats per lletres majúscules, com
per exemple $P(x)$ o $Q(m,n)$, afegint les variables que estan presents
entre parèntesis. Per exemple, l'enunciat \enquote{$n$ és un nombre primer}
és un predicat que representem per $P(n)$, i l'enunciat \enquote{$m$ és més gran o igual que $n$}, o abreujadament, \enquote{$m\geq n$} per $Q(m,n)$.
Llavors, $P(2)$ és la proposició \enquote{$2$ és un nombre primer} que
és vertadera, i $Q(3,4)$ és la proposició \enquote{$3\geq4$} que és
falsa.

Dels cinc últims enunciats anteriors que no són proposicions els
quatre primers són predicats. El primer és el predicat \enquote{ser un
nombre parell}, que si el denotem per $P$, aleshores l'enunciat és $P(x)$. El segon és el predicat \enquote{ser una recta paral·lela a
$x+y=1$}, que si el simbolitzem per $Q$ i escrivim una recta del pla com $ax+by=c$, aleshores l'enunciat és $Q(a,b,c)$; per exemple, $Q(2,2,-1)$
és la proposició \enquote{$2x+2y=-1$ és una recta
paral·lela a $x+y=1$} que és vertadera. De fet el
tercer és una família de predicats que depèn del paràmetre $a$ que és un nombre irracional. En concret, si el designem per $P_{a}$ i
considerem que la variable $x$ pren valors a $\mathbb{Z}$, aleshores
l'enunciat és $P_{a}(x)$; per exemple, $P_{\sqrt{2}}(4)$ és la
proposició \enquote{$4$ és un nombre enter parell més gran que el nombre
irracional $\sqrt{2}$} que és vertadera. Finalment, l'últim és el
predicat \enquote{$f$ és la funció inversa de $g$}, que si el denotem per $S$
i $f,g$ són variables que denoten funcions, aleshores l'enunciat és $S(f,g)$. Per exemple, si $f:\mathbb{R}\to\mathbb{R}$ i $g:\mathbb{R}\to\mathbb{R}$ són les funcions definides per $f(t)=\sqrt[3]{t}$ i $g(t)=t^{3}$, aleshores $S(f,g)$ és la
proposició \enquote{$f$ és la funció inversa de $g$} que és
vertadera. En
canvi, l'últim no és un predicat perquè al substituir la
variable $x$ per un número s'obté un altre número, però en
cap cas una proposició; diem que és una expressió algebraica.
Aquestes expressions estan compostes per variables, constants numèriques
i els símbols de les quatre operacions fonamentals de l'aritmètica.


Per acabar aquesta secció volem advertir que convé distingir entre una proposició matemàtica i una afirmació que parla sobre aquesta proposició. Per exemple, si volem dir que la proposició \enquote{$3\geq 4$} és falsa, no escrivim simplement
\begin{equation*}
3\geq 4\text{ és falsa},
\end{equation*}
perquè aquí barregem l'enunciat matemàtic amb el comentari que fem sobre el seu valor de veritat. És millor escriure, en llenguatge ordinari,
\begin{equation*}
\text{la proposició \enquote{$3\geq 4$} és falsa.}
\end{equation*}
Aquesta distinció no cal portar-la fins a l'extrem en els càlculs quotidians, però és important tenir-la clara quan parlem de veritat, falsedat, demostració o contradicció.

\begin{observacio}
En aquesta unitat adoptem un punt de vista elemental: treballem dins del context matemàtic i, mentre no es digui el contrari, quan parlem de proposicions ens referim a proposicions matemàtiques.

Si $p$ és una proposició matemàtica determinada, l'afirmació \enquote{$p$ és falsa} no és una nova proposició matemàtica del mateix tipus que $p$: no parla directament de nombres, figures o funcions, sinó de la proposició $p$ i del seu valor de veritat.

Així, si $p$ és falsa, l'afirmació \enquote{$p$ és falsa} és vertadera; i, si $p$ és vertadera, l'afirmació \enquote{$p$ és falsa} és falsa. Per exemple, si $p$ denota la proposició \enquote{$3\geq4$}, aleshores l'afirmació \enquote{$p$ és falsa} és vertadera.

Quan una afirmació parla sobre ella mateixa, poden aparèixer dificultats lògiques: un raonament que sembla correcte pot conduir a una contradicció. Més endavant, aquest tipus de situació rep el nom de paradoxa. L'exemple clàssic és la paradoxa del mentider: \enquote{aquesta frase és falsa}. Si fos vertadera, aleshores seria falsa; si fos falsa, aleshores seria vertadera. Aquesta mena d'autoreferència és la que cal evitar.

Una situació semblant apareix en teoria de conjunts amb la paradoxa de Russell. Si intentem formar el conjunt
\begin{equation*}
A=\left\{x:x\notin x\right\},
\end{equation*}
aleshores preguntar si $A\in A$ porta a una contradicció: si $A\in A$, per la definició de $A$ hauríem de tenir $A\notin A$; i si $A\notin A$, aleshores $A$ compleix la propietat que defineix $A$, i per tant $A\in A$.

Per evitar aquests problemes, en matemàtiques és habitual separar el llenguatge que fem servir per escriure els enunciats matemàtics del llenguatge amb què parlem sobre aquests enunciats.
\end{observacio}

\section{Connectives lògiques}

En teoria de conjunts combinem o modifiquem els conjunts amb operacions com
la unió o la intersecció. De la mateixa manera, en aritmètica
modifiquem nombres amb operacions com la suma o la multiplicació. Això també passa en lògica. Tenim algunes operacions per combinar o
modificar proposicions o predicats; aquestes operacions es tradueixen en el
llenguatge de les matemàtiques per les paraules \enquote{i}, \enquote{o}, \enquote{no}, \enquote{si ...
llavors ...} i \enquote{... si i només si ... }. És habitual simbolitzar
aquestes operacions amb els símbols característics del llenguatge
de la lògica. Aquí ho farem per una qüestió d'economia sintàctica. Cal tenir present que a matemàtiques, aquestes paraules
tenen significats precisos, que donarem a l'apartat següent. En alguns
casos, el significat matemàtic d'aquestes paraules difereix lleugerament
o són més precisos que l'ús comú que fem en la vida
quotidiana.

En general, la regla sintàctica és que a partir de dues proposicions
$p$ i $q$ podem formar altres proposicions mitjançant les operacions lògiques, també anomenades \textbf{connectives lògiques}. La
regla semàntica és que la veritat o falsedat d'aquestes noves
proposicions dependrà de la veritat o falsedat de $p$ i $q$, i serà
definida per les anomenades \textbf{taules de veritat}. A una proposició
se l'assigna el valor de veritat que simbolitzem per V quan la proposició
és vertadera, i F quan és falsa.

\subsection{Conjunció}

Si $p$ i $q$ són dues proposicions, llavors la \textbf{conjunció}
d'aquestes dues proposicions és la proposició \enquote{$p$ i $q$},
simbolitzada per $p\wedge q$, i que es defineix de la següent manera:

\begin{itemize}
\item És vertadera, quan $p$ i $q$ són ambdues vertaderes;

\item És falsa, quan $p$ és falsa o $q$ és falsa o ambdues són falses.
\end{itemize}

Per exemple, la proposició \enquote{$1<\sqrt{2}<2$} és la conjunció de
les proposicions \enquote{$\sqrt{2}>1$} i \enquote{$\sqrt{2}<2$}, que són totes dues
vertaderes i, per tant, és vertadera.

D'acord amb la definició anterior, s'obté la taula de veritat de la
conjunció lògica:\begin{equation*}
\begin{tabular}{ll|l}
$p$ & $q$ & $p\wedge q$ \\ \hline
\multicolumn{1}{c}{V} & \multicolumn{1}{c|}{V} & \multicolumn{1}{|c}{V} \\
\multicolumn{1}{c}{V} & \multicolumn{1}{c|}{F} & \multicolumn{1}{|c}{F} \\
\multicolumn{1}{c}{F} & \multicolumn{1}{c|}{V} & \multicolumn{1}{|c}{F} \\
\multicolumn{1}{c}{F} & \multicolumn{1}{c|}{F} & \multicolumn{1}{|c}{F}\end{tabular}
\text{.}
\end{equation*}
Observem que la conjunció de dues proposicions només és
vertadera quan les dues proposicions són vertaderes.

\subsection{Disjunció}

Si $p$ i $q$ són dues proposicions, llavors la \textbf{disjunció}
d'aquestes dues proposicions és la proposició \enquote{$p$ o $q$},
simbolitzada per $p\vee q$, i que es defineix de la següent manera:

\begin{itemize}
\item És vertadera, quan almenys una de les dues proposicions és
vertadera;

\item És falsa, quan $p$ i $q$ són ambdues falses.
\end{itemize}

Per exemple, la proposició \enquote{$\pi\in\left( -\infty,-1\right) \cup\left(
1,+\infty\right) $} és la disjunció de les proposicions \enquote{$\pi<-1$} o
\enquote{$\pi>1$}, la primera és falsa i la segona és vertadera i, per tant,
és vertadera. La taula de veritat d'aquesta connectiva és\begin{equation*}
\begin{tabular}{ll|l}
$p$ & $q$ & $p\vee q$ \\ \hline
\multicolumn{1}{c}{V} & \multicolumn{1}{c|}{V} & \multicolumn{1}{|c}{V} \\
\multicolumn{1}{c}{V} & \multicolumn{1}{c|}{F} & \multicolumn{1}{|c}{V} \\
\multicolumn{1}{c}{F} & \multicolumn{1}{c|}{V} & \multicolumn{1}{|c}{V} \\
\multicolumn{1}{c}{F} & \multicolumn{1}{c|}{F} & \multicolumn{1}{|c}{F}\end{tabular}
\text{.}
\end{equation*}
Observem que la disjunció de dues proposicions és falsa quan les
dues proposicions són falses. Destaquem l'ús \textbf{inclusiu} que
fem d'aquesta connectiva respecte de l'exclusiu que fem servir habitualment
a la vida quotidiana i que fa que la proposició sigui vertadera quan només una de les proposicions és vertadera.

\subsection{Negació}

Si $p$ és una proposició, llavors la \textbf{negació} d'aquesta
proposició és la proposició \enquote{no $p$}, simbolitzada per $\lnot p$, que es defineix de la següent manera:

\begin{itemize}
\item És vertadera, quan $p$ és falsa;

\item És falsa, quan $p$ és vertadera.
\end{itemize}

Per exemple, la proposició \enquote{$\sqrt{3}\notin\left( 0,1\right) $} és
la negació de la proposició \enquote{$\sqrt{3}\in\left( 0,1\right) $}, que
és falsa i, per tant, és vertadera. La taula de veritat d'aquesta
connectiva és\begin{equation*}
\begin{tabular}{l|l}
$p$ & $\lnot p$ \\ \hline
\multicolumn{1}{c|}{V} & \multicolumn{1}{|c}{F} \\
\multicolumn{1}{c|}{F} & \multicolumn{1}{|c}{V}\end{tabular}
\text{.}
\end{equation*}

\subsection{Condicional}

Si $p$ i $q$ són dues proposicions, llavors el \textbf{condicional}
d'aquestes dues proposicions és la proposició \enquote{si $p$, llavors $q$},
simbolitzada per $p\longrightarrow q$, i que es defineix de la següent
manera:

\begin{itemize}
\item És vertadera, quan $p$ i $q$ són ambdues vertaderes o $p$ és falsa;

\item És falsa, quan $p$ és vertadera i $q$ és falsa.
\end{itemize}

Una proposició condicional \enquote{si $p$, llavors $q$} també se'n diu
\textbf{implicació} i diem que \enquote{$p$ implica $q$}, i que $p$ és l'\textbf{antecedent} i $q$ és el \textbf{conseqüent} de la implicació. La interpretació que fem d'aquesta connectiva sorprèn una
mica en el seu ús doncs, per exemple, en l'enunciat \enquote{Si hi ha dues
rectes en el pla que són paral·leles, llavors 2 és un nombre primer} no hi ha cap relació entre l'antecedent \enquote{hi ha dues
rectes en el pla que són paral·leles} i el conseqüent \enquote{$2$ és un nombre primer}, i, és matemàticament correcte,
independentment de si l'antecedent és vertader o fals.

Segons la definició que hem donat, la taula de veritat d'aquesta
connectiva és\begin{equation*}
\begin{tabular}{ll|l}
$p$ & $q$ & $p\longrightarrow q$ \\ \hline
\multicolumn{1}{c}{V} & \multicolumn{1}{c|}{V} & \multicolumn{1}{|c}{V} \\
\multicolumn{1}{c}{V} & \multicolumn{1}{c|}{F} & \multicolumn{1}{|c}{F} \\
\multicolumn{1}{c}{F} & \multicolumn{1}{c|}{V} & \multicolumn{1}{|c}{V} \\
\multicolumn{1}{c}{F} & \multicolumn{1}{c|}{F} & \multicolumn{1}{|c}{V}\end{tabular}
\text{.}
\end{equation*}

Observem que només és falsa quan l'antecedent és vertader i el
conseqüent és fals. El condicional de dues proposicions $p$ i $q$
podem trobar-lo en proposicions escrites al català de moltes maneres:
(1) $q$ si $p$; (2) si $p$, $q$; (3) $p$ només si $q$; o (4) $q$ sempre
que $p$.

Per exemple, la proposició \enquote{700 és divisible per 4 perquè el
nombre acaba amb dos zeros} és el condicional \enquote{$P(700)$ implica $Q(700)$}, on $P$ és el predicat \enquote{tenir zero a les decenes i també a les
unitats} i $Q$ és \enquote{ser divisible per 4}.

\subsection{Bicondicional}

Si $p$ i $q$ són dues proposicions, llavors el \textbf{bicondicional}
d'aquestes dues proposicions és la proposició \enquote{$p$ si i només si
$q$}, simbolitzada per $p\longleftrightarrow q$, i que es defineix de la següent manera:

\begin{itemize}
\item És vertadera, quan $p$ i $q$ són ambdues vertaderes o falses;

\item És falsa, quan $p$ és vertadera i $q$ és falsa, o $p$ és falsa i $q$ és vertadera.
\end{itemize}

Una proposició bicondicional \enquote{$p$ si i només si $q$} també se'n
diu \textbf{doble implicació} i diem que \enquote{$p$ implica $q$ i $q$ implica $p$}. La taula de veritat d'aquesta connectiva és\begin{equation*}
\begin{tabular}{ll|l}
$p$ & $q$ & $p\longleftrightarrow q$ \\ \hline
\multicolumn{1}{c}{V} & \multicolumn{1}{c|}{V} & \multicolumn{1}{|c}{V} \\
\multicolumn{1}{c}{V} & \multicolumn{1}{c|}{F} & \multicolumn{1}{|c}{F} \\
\multicolumn{1}{c}{F} & \multicolumn{1}{c|}{V} & \multicolumn{1}{|c}{F} \\
\multicolumn{1}{c}{F} & \multicolumn{1}{c|}{F} & \multicolumn{1}{|c}{V}\end{tabular}
\text{.}
\end{equation*}

Per exemple, la proposició \enquote{98 és múltiple de 6 si i només
si 98 és divisible per 2 i també per 3} és el bicondicional \enquote{$P(98)$ si i només si $Q(2)$ i $R(3)$}, on $P$ és el predicat \enquote{ser múltiple de 6}, $Q$ és \enquote{ser divisible per 2} i $R$ és \enquote{ser divisible
per 3}.

\begin{exemple}
Escriu fent ús de les connectives les proposicions següents: (1) $-4\leq-1$; (2) $\left\vert -3\right\vert =\left\vert 2\right\vert $; (3) $-1\in\left[ 0,1\right] $; (4) $2\notin\left[ 0,1\right] $; (5) $\sqrt {2}<2<2\sqrt{2}$ si $1<\sqrt{2}<2$; (6) $\sqrt[3]{1331}=11$ si i només si $11^{3}=1331$.
\end{exemple}

\begin{solucio}
(1) La proposició $-4\leq-1$ es pot expressar com la disjunció: \enquote{$-4<-1$} o \enquote{$-4=-1$}. (2) La proposició $\left\vert -3\right\vert
=\left\vert 2\right\vert $ és la disjunció: \enquote{$-3=2$} o'$-3=-2$'. (3)
La proposició $-1\in\left[ 0,1\right] $ és la conjunció: \enquote{$-1\geq0$} i \enquote{$-1\leq1$}. (4) La proposició $2\notin\left[ 0,1\right] $
és la negació: \enquote{no $2\in\left[ 0,1\right] $}. (5) La proposició \enquote{$\sqrt{2}<2<2\sqrt{2}$ si $1<\sqrt{2}<2$} és el condicional \enquote{$1<\sqrt {2}<2$ implica $\sqrt{2}<2<2\sqrt{2}$}. (6) La proposició \enquote{$\sqrt[3]{1331}=11$ si i només si $11^{3}=1331$} és el bicondicional \enquote{ $11^{3}=1331$
implica $\sqrt[3]{1331}=11$ i $\sqrt[3]{1331}=11$ implica $11^{3}=1331$}.
\end{solucio}

\section{Proposicions elementals i compostes}

Proposicions \textbf{elementals} són aquelles que no posseeixen cap
operador lògic. Les proposicions \textbf{compostes} estan formades per
altres proposicions i operadors lògics. Per exemple, a partir de dues
proposicions elementals $p$ i $q$, podem formar com hem vist $\lnot q$ i $p\vee q$. Però també $p\wedge\lnot q$, i finalment, $\left( p\vee
q\right) \longrightarrow\left( p\wedge\lnot q\right).$ Hem utilitzat els parèntesis per evitar ambigüetats. De fet, també hem utilitzat la
convenció estàndard segons el qual el símbol de negació $\lnot$ té prioritat sobre les altres quatre connectives, però cap
d'aquestes té prioritat sobre les altres. Per altra banda, és clar
que $p\wedge\longrightarrow\lnot q$ no té cap sentit perquè no està ben formada. La raó, per exemple, està en què $\wedge$
connecta dues proposicions i, en el nostre cas, $\longrightarrow\lnot q$ no
ho és.

Volem ara formalitzar l'enunciat següent: \enquote{Si no és cert que 34043
sigui una suma de dos quadrats, llavors 34043 és senar o és
divisible per 3}. Identifiquem les següents proposicions simples: $p=$
\enquote{34043 sigui una suma de dos quadrats}, $q=$ \enquote{34043 és senar}, i $r=$
\enquote{34043 és divisible per 3}. Llavors l'enunciat s'escriu simbòlicament d'aquesta manera tenint en compte les connectives lògiques: $\lnot p\longrightarrow\left( q\vee r\right).$ Un altre exemple és la proposició \enquote{273 és divisible per 91, si 273 és múltiple de 7 i múltiple de 13} és el condicional: \enquote{$P(273)\longrightarrow\left(
Q(273)\wedge R(273)\right) $}, on $P$ és el predicat \enquote{ser divisible per
91}, $Q$ és \enquote{ser mútiple de 7} i $R$ és \enquote{ser múltiple de 13}.

\bigskip

Donada una proposició composta, volem determinar el seu valor de veritat
coneixent el valor de veritat de les proposicions simples que la conformen.
Suposem la interpretació segons la qual les proposicions simples $p$ i $q $ són vertaderes i $r$ i $s$ són falses. Llavors, quin és el
significat de la proposició composta $\lnot\left( p\vee q\right)
\longrightarrow\left( r\wedge\lnot s\right) $?

\begin{equation*}
\begin{tabular}{cccccccc|c}
$p$ & $q$ & $r$ & $s$ & $\lnot s$ & $p\vee q$ & $\alpha=\lnot\left( p\vee
q\right) $ & $\beta=\left( r\wedge\lnot s\right) $ & $\alpha\longrightarrow
\beta$ \\ \hline
V & V & F & F & V & V & F & F & V\end{tabular}
\end{equation*}

De la taula deduïm que la proposició és vertadera.

\bigskip

Donat el valor de veritat d'una proposició composta, ara volem
determinar el valor de veritat de les proposicions simples que la conformen.
Suposem que la proposició composta $\left( p\wedge\lnot q\right)
\longrightarrow r$ és falsa. Llavors, com la connectiva principal
d'aquesta proposició és el condicional. Atès que aquesta
implicació té un valor de veritat fals únicament quan
l'antecedent és vertader i el conseqüent és fals, s'obté que
$p\wedge\lnot q$ és vertadera, i $r$ ha de ser falsa. Ara bé $p\wedge\lnot q$ és vertadera només si $p$ és vertadera i $q$
és falsa. Per tant, $p$ és vertadera i $q$ i $r$ són falses.

\section{Formes proposicionals}

Com hem fet fins ara les lletres $p,q,r,...$ s'usen com a variables que
representan a proposicions. El valor de veritat de $p$, per exemple, és
desconegut mentre no s'especifica la seva interpretació, és a dir,
no és coneix el significat de la proposició que representa.

Anomenem \textbf{formes proposicionals} a les expressions formals ben
formades constituïdes per variables proposicionals i les connectives lògiques que les relacionen. Per exemple, $A(p,q,r)=\left( \left( p\wedge
q\right) \longrightarrow\left( r\vee\lnot p\right) \right) \wedge r$ és
una forma proposional que conté tres variables $p,q$ i $r$. Les formes
proposicionals no són proposicions, però si cada variable
proposicional és reemplaçada per una proposició simple o
composta, la forma proposicional es converteix en una proposició. Si
reemplacem les variables proposicionals per proposicions vertaderes o
falses, el nombre de proposicions que es generen és $2^{n}$, sent $n$ el
nombre de variables proposicionals. A la taula de veritat que es genera se
l'anomena a vegades \textbf{matriu} de la forma proposicional. La matriu de
la forma $A(p,q,r)$ té $2^{3}=8$ interpretacions posibles:\begin{equation*}
\begin{tabular}{cc|c|c|c|ccc}
$p$ & $q$ & $r$ & $\lnot p$ & $\alpha=p\wedge q$ & $\beta=r\vee\lnot p$ & $\alpha\longrightarrow\beta$ & $\left( \alpha\longrightarrow\beta\right)
\wedge r$ \\ \hline
V & V & V & F & V & V & V & V \\
V & V & F & F & V & F & F & F \\
V & F & V & F & F & V & V & V \\
V & F & F & F & F & F & V & F \\
F & V & V & V & F & V & V & V \\
F & V & F & V & F & V & V & F \\
F & F & V & V & F & V & V & V \\
F & F & F & V & F & V & V & F\end{tabular}
\
\end{equation*}
Per exemple, la forma es converteix en una proposició vertadera quan les
variables proposicionals $p,q$ i $r$ s'interpreten respectivament com a
proposicions que són vertadera, falsa i vertadera (fila 3 de la matriu).

\begin{exemple}
Construeix la matriu de la forma $A(p,q)=\left( p\vee q\right)
\longrightarrow\left( p\wedge\lnot q\right).$ Quina interpretació
permet assegurar que la forma es converteix en una proposició vertadera?
\end{exemple}

\begin{solucio}
La forma té dues variables proposicionals i, per tant, $2^{2}=4$
interpretacions possibles. La matriu de la forma és:
\begin{equation*}
\begin{tabular}{cc|c|c|cc}
$p$ & $q$ & $\lnot q$ & $p\vee q$ & $p\wedge\lnot q$ & $\left( p\vee
q\right) \longrightarrow\left( p\wedge\lnot q\right) $ \\ \hline
V & V & F & V & F & F \\
V & F & V & V & V & V \\
F & V & F & V & F & F \\
F & F & V & F & F & V\end{tabular}
\
\end{equation*}
Observem que una proposició d'aquesta forma només serà vertadera
quan substituïm $q$ per una proposició falsa.
\end{solucio}

\section{Tautologies i contradiccions}

Una \textbf{tautologia} és una proposició que és vertadera per
necessitat lògica. Una tautologia característica de la lògica clàssica és la proposició \enquote{$p$ o no $p$}, sent $p$ qualsevol
proposició. Aquesta tautologia se la coneix també amb el nom de
\textbf{principi del tercer exclòs}, segons el qual la disjunció
d'una proposició i la seva negació és sempre vertadera. En
efecte, la taula de veritat d'aquesta proposició és\begin{equation*}
\begin{tabular}{ll|l}
$p$ & $\lnot p$ & $p\vee\lnot p$ \\ \hline
\multicolumn{1}{c}{V} & \multicolumn{1}{c|}{F} & \multicolumn{1}{|c}{V} \\
\multicolumn{1}{c}{F} & \multicolumn{1}{c|}{V} & \multicolumn{1}{|c}{V}\end{tabular}
\text{,}
\end{equation*}
on veiem que l'última columna hi ha només el valor V com era
d'esperar. Es evident que la tautologia no té cap interés matemàtic perquè no informa de res, però permet simplificar a vegades les
demostracions com veurem més endavant.

\bigskip

Ara bé, el concepte de tautologia té més sentit quan en lloc de
tractar amb proposicions treballem en formes. Una forma és tautologia
quan la matriu de la forma només genera proposicions vertaderes; en unes
altres paraules, quan és certa en totes les circumstàncies
possibles. Per exemple, la forma $A(p,q)=(p\wedge q)\longrightarrow(p\vee q)
$ és tautologia perquè la seva taula és:\begin{equation*}
\begin{tabular}{ll|l|l|l}
$p$ & $q$ & $p\wedge q$ & $p\vee q$ & $(p\wedge q)\longrightarrow(p\vee q)$
\\ \hline
\multicolumn{1}{c}{V} & \multicolumn{1}{c|}{V} & \multicolumn{1}{|c|}{V} &
\multicolumn{1}{|c|}{V} & \multicolumn{1}{|c}{V} \\
\multicolumn{1}{c}{V} & \multicolumn{1}{c|}{F} & \multicolumn{1}{|c|}{F} &
\multicolumn{1}{|c|}{V} & \multicolumn{1}{|c}{V} \\
\multicolumn{1}{c}{F} & \multicolumn{1}{c|}{V} & \multicolumn{1}{|c|}{F} &
\multicolumn{1}{|c|}{V} & \multicolumn{1}{|c}{V} \\
\multicolumn{1}{c}{F} & \multicolumn{1}{c|}{F} & \multicolumn{1}{|c|}{F} &
\multicolumn{1}{|c|}{F} & \multicolumn{1}{|c}{V}\end{tabular}
\end{equation*}
D'aquí surt el fet següent: qualsevol proposició composta de la
forma \enquote{$(p\wedge q)\longrightarrow(p\vee q)$} és una tautologia,
independentment del que siguin les proposicions $p$ i $q$.

\bigskip

Una \textbf{contradicció} és una proposició que és falsa per
necessitat lògica. De fet, tota contradicció és la negació
d'una tautologia i, al contrari, tota tautologia és la negació d'una
contradicció. La negació del principi del tercer exclòs és
una contradicció coneguda com el \textbf{principi de contradicció},
segons el qual la conjunció d'una proposició i la seva negació
és una proposició sempre falsa. És la proposició \enquote{$p$ i no$~p $}, sent $p$ qualsevol proposició. La taula de veritat d'aquesta
proposició és\begin{equation*}
\begin{tabular}{ll|l}
$p$ & $\lnot p$ & $p\wedge\lnot p$ \\ \hline
\multicolumn{1}{c}{V} & \multicolumn{1}{c|}{F} & \multicolumn{1}{|c}{F} \\
\multicolumn{1}{c}{F} & \multicolumn{1}{c|}{V} & \multicolumn{1}{|c}{F}\end{tabular}
\text{,}
\end{equation*}
i veiem que l'última columna només té el valor F. De la mateixa
manera que la tautologia, la contradicció té més sentit quan
s'aplica a formes proposicionals. Una forma és contradicció quan la
seva matriu només genera proposicions falses. Per exemple, segons el
principi de contradicció la forma $\left( p\longrightarrow q\right)
\wedge\lnot\left( p\longrightarrow q\right) $ és una contradicció.
Podem comprovar-ho fent la seva matriu:\begin{equation*}
\begin{tabular}{ll|l|l|l}
$p$ & $q$ & $p\longrightarrow q$ & $\lnot\left( p\longrightarrow q\right) $
& $(p\longrightarrow q)\wedge\lnot\left( p\longrightarrow q\right) $ \\
\hline
\multicolumn{1}{c}{V} & \multicolumn{1}{c|}{V} & \multicolumn{1}{|c|}{V} &
\multicolumn{1}{|c|}{F} & \multicolumn{1}{|c}{F} \\
\multicolumn{1}{c}{V} & \multicolumn{1}{c|}{F} & \multicolumn{1}{|c|}{F} &
\multicolumn{1}{|c|}{V} & \multicolumn{1}{|c}{F} \\
\multicolumn{1}{c}{F} & \multicolumn{1}{c|}{V} & \multicolumn{1}{|c|}{V} &
\multicolumn{1}{|c|}{F} & \multicolumn{1}{|c}{F} \\
\multicolumn{1}{c}{F} & \multicolumn{1}{c|}{F} & \multicolumn{1}{|c|}{V} &
\multicolumn{1}{|c|}{F} & \multicolumn{1}{|c}{F}\end{tabular}
\text{.}
\end{equation*}

Per acabar, quan una forma proposicional no és tautologia ni contradicció es diu que és \textbf{contingència}; en unes altres \
paraules, quan admet algunes proposicions vertaderes i altres falses per als
valors de veritat de les variables proposicionals. Un exemple de contingència és la forma $(p\vee q)\longrightarrow(p\wedge q)$ perquè la
seva matriu és:\begin{equation*}
\begin{tabular}{cc|c|c|c|}
$p$ & $q$ & $p\vee q$ & $p\wedge q$ & $(p\vee q)\longrightarrow(p\wedge q)$
\\ \hline
V & V & V & V & V \\
V & F & V & F & F \\
F & V & V & F & F \\
F & F & F & F & V\end{tabular}
\text{.}
\end{equation*}

\bigskip

Les formes proposicionals poden ser connectades amb operadors lògics per
a formar noves formes proposicionals. D'aquesta manera, si $A$ i $B$ són
formes, llavors $\lnot A$, $A\wedge B$, $A\vee B$, $A\longrightarrow B$ i $A\longleftrightarrow B$ representen noves formes proposicionals. Per
exemple, si $A(p,q)=\lnot p\longrightarrow q$ i $B(p,r,s)=s\longrightarrow\left( p\vee\lnot r\right) $, aleshores $C(p,q,r,s)=A(p,q)\longrightarrow
B(p,r,s)$, és a dir, $C(p,q,r,s)=\left( \lnot p\longrightarrow q\right)
\longrightarrow\left( s\longrightarrow\left( p\vee\lnot r\right) \right).$
Observem que $A$ té $2^{2}=4$ interpretacions, $B$ en té $2^{3}=8$
interpretacions, i $C$ en té $2^{4}=16$ interpretacions.

\begin{exemple}
Si $p,q$ i $r$ són variables proposicionals, llavors
\[
\left((p\longrightarrow q)\wedge(q\longrightarrow r)\right)
\longrightarrow(p\longrightarrow r)
\]
és una forma proposicional tautològica?
\end{exemple}

\begin{solucio}
Sí. Aquesta forma expressa la transitivitat del condicional: si $p$ implica
$q$ i $q$ implica $r$, aleshores $p$ implica $r$. Ho podem comprovar amb la
taula de veritat:
\begin{equation*}
\resizebox{\textwidth}{!}{\begin{tabular}{cccccc}
$p$ & $q$ & $r$ & $p\longrightarrow q$ & $q\longrightarrow r$ &
$\left((p\longrightarrow q)\wedge(q\longrightarrow r)\right)\longrightarrow(p\longrightarrow r)$ \\ \hline
V & V & V & V & V & V \\
V & V & F & V & F & V \\
V & F & V & F & V & V \\
V & F & F & F & V & V \\
F & V & V & V & V & V \\
F & V & F & V & F & V \\
F & F & V & V & V & V \\
F & F & F & V & V & V\end{tabular}}
\end{equation*}
Com que l'última columna només conté valors V, la forma és una tautologia.
\end{solucio}

\section{Implicació lògica\label{sec:implicacio-logica}}

Si $A$ i $B$ són dues formes proposicionals, llavors es diu que $A$
\textbf{implica lògicament} $B$ quan la forma $A\longrightarrow B$ és tautologia. Cal destacar que la implicació lògica és una relació entre formes proposicionals, que denotem per $\Longrightarrow$.
D'aquesta manera, quan escrivim $A\Longrightarrow B$, que llegim
abreujadament com \enquote{$A$ implica $B$}, significa que $A\longrightarrow B$ és tautologia. Observa que hem usat en aquest document la paraula \enquote{implica}
de dues maneres diferents. La primera, al parlar del condicional de dues
proposicions $p$ i $q$, $p\rightarrow q$, què no necessàriament es
una tautologia, mentre que la segona, fa un moment, quan la implicació sí és necessàriament una tautologia.

A la pràctica, el fet de tenir la implicació lògica $A\Longrightarrow B$ permet fer el següent argument: si $A\Longrightarrow
B$ i $A$ és tautologia, aleshores necessariàment $B$ és
tautologia. D'aquí surt una interpretació molt comú en matemàtiques: quan $A\Longrightarrow B$, aleshores $B$ es diu que és
\textbf{condició necessària} per $A$, i també, que $A$ és
\textbf{condició suficient} per $B$.

Per exemple, suposem que $n$ és un enter positiu, aleshores l'enunciat
\enquote{si $n$ és divisible per $4$, llavors $n$ és divisible per $2$},
és cert i pot escriure's com $P(n)\longrightarrow Q(n)$, on $P=$ \enquote{és
divisible per $4$} i $Q=$ \enquote{és divisible per $2$}. La condició \enquote{$n$
és divisible per $4$} és suficient perquè \enquote{$n$ sigui divisible
per $2$}, és a dir, només cal que $n$ sigui divisible per $4$ perquè sigui divisible per $2$. En canvi, la condició \enquote{$n$ és
divisible per $2$} és necessària perquè \enquote{$n$ és divisible
per $4$}, o sigui, si $n$ no fos divisible per $2$, aleshores no seria
divisible per 4.

Més endavant veurem que les implicacions lògiques seran
extremadament útils per construir raonaments vàlids. En particular,
s'utilitzaran àmpliament les següents implicacions lògiques: Si $A,B$ i $C$ són formes proposicionals, aleshores

\begin{enumerate}
\item $(A\longrightarrow B)\wedge A$ $\Longrightarrow B$

\item $(A\longrightarrow B)\wedge\lnot B\Longrightarrow\lnot A$

\item (i) $A\wedge B$ $\Longrightarrow A$; (ii) $A\wedge B$ $\Longrightarrow
B$

\item (i) $A\Longrightarrow A\vee B$; (ii) $B\Longrightarrow A\vee B$

\item (i) $\left( A\vee B\right) \wedge\lnot B\Longrightarrow A$; (ii) $\left( A\vee B\right) \wedge\lnot A\Longrightarrow B$

\item (i) $A\longleftrightarrow B\Longrightarrow A\longrightarrow B$; (ii) $A\longleftrightarrow B\Longrightarrow B\longrightarrow A$

\item $\left( A\longrightarrow B\right) \wedge\left( B\longrightarrow
A\right) \Longrightarrow A\longleftrightarrow B$

\item $\left( A\longrightarrow B\right) \wedge\left( B\longrightarrow
C\right) \Longrightarrow A\longrightarrow C$
\end{enumerate}

Per veure que $(A\longrightarrow B)\wedge A$ $\Longrightarrow B$ hem de
comprovar que $(\left( A\longrightarrow B)\wedge A\right) $ $\longrightarrow
B$ és tautologia. Construïm la taula de veritat corresponent:\begin{equation*}
\begin{tabular}{lllll}
$A$ & $B$ & $A\longrightarrow B$ & $\left( A\longrightarrow B\right) \wedge
A $ & $\left( \left( A\longrightarrow B\right) \wedge A\right)
\longrightarrow B$ \\ \hline
\multicolumn{1}{c}{V} & \multicolumn{1}{c|}{V} & \multicolumn{1}{c}{V} &
\multicolumn{1}{c}{V} & \multicolumn{1}{c}{V} \\
\multicolumn{1}{c}{V} & \multicolumn{1}{c|}{F} & \multicolumn{1}{c}{F} &
\multicolumn{1}{c}{F} & \multicolumn{1}{c}{V} \\
\multicolumn{1}{c}{F} & \multicolumn{1}{c|}{V} & \multicolumn{1}{c}{V} &
\multicolumn{1}{c}{F} & \multicolumn{1}{c}{V} \\
\multicolumn{1}{c}{F} & \multicolumn{1}{c|}{F} & \multicolumn{1}{c}{V} &
\multicolumn{1}{c}{F} & \multicolumn{1}{c}{V}\end{tabular}
\
\end{equation*}
i observem que efectivament és tautologia. Anàlogament es fa per les
altres.

\section{Equivalència lògica\label{sec:equivalencia-logica}}

Si $A$ i $B$ són dues formes proposicionals, llavors es diu que $A$ i $B$
són \textbf{equivalents lògicament }quan la forma $A\longleftrightarrow B$ és tautologia. Com la implicació lògica,
l'equivalència és una relació entre formes proposicionals que
denotem per $\Longleftrightarrow$. D'aquesta manera, quan escrivim $A\Longleftrightarrow B$, que llegim \enquote{$A$ és equivalent a $B$} o \enquote{$A$ i $B$ són equivalents}, significa que $A\longleftrightarrow B$ és
tautologia.

A la pràctica, el fet de tenir l'equivalència $A\Longleftrightarrow
B $ permet fer el següent argument: Si $A\Longleftrightarrow B$, i $A$
(resp. $B$) és tautologia, aleshores necessàriament $B$ (resp. $A$)
és tautologia, i també, d'aquí surt una interpretació molt
comuna en matemàtiques: quan $A\Longleftrightarrow B$, aleshores es diu
que $A$ és \textbf{condició necessària i suficient} per $B$ o a
l'inrevés. De fet, això es dedueix de l'equivalència següent: $A\longleftrightarrow B\Longleftrightarrow\left( A\longrightarrow
B\right) \wedge\left( B\longrightarrow A\right) $ perquè

\begin{equation*}
\begin{tabular}{lllllll}
$A$ & $B$ & $\alpha=A\longrightarrow B$ & $\beta=B\longrightarrow A$ & $\gamma=A\longleftrightarrow B$ & $\alpha\wedge\beta$ & $\gamma
\longleftrightarrow\left( \alpha\wedge\beta\right) $ \\ \hline
\multicolumn{1}{c}{V} & \multicolumn{1}{c|}{V} & \multicolumn{1}{c}{V} &
\multicolumn{1}{c}{V} & \multicolumn{1}{c}{V} & \multicolumn{1}{c}{V} &
\multicolumn{1}{c}{V} \\
\multicolumn{1}{c}{V} & \multicolumn{1}{c|}{F} & \multicolumn{1}{c}{F} &
\multicolumn{1}{c}{V} & \multicolumn{1}{c}{F} & \multicolumn{1}{c}{F} &
\multicolumn{1}{c}{V} \\
\multicolumn{1}{c}{F} & \multicolumn{1}{c|}{V} & \multicolumn{1}{c}{V} &
\multicolumn{1}{c}{F} & \multicolumn{1}{c}{F} & \multicolumn{1}{c}{F} &
\multicolumn{1}{c}{V} \\
\multicolumn{1}{c}{F} & \multicolumn{1}{c|}{F} & \multicolumn{1}{c}{V} &
\multicolumn{1}{c}{V} & \multicolumn{1}{c}{V} & \multicolumn{1}{c}{V} &
\multicolumn{1}{c}{V}\end{tabular}
\end{equation*}

A continuació s'enumeren algunes equivalències lògiques que
seran particularment útils i al seu costat el nom pel qual són
conegudes:

\begin{enumerate}
\item $\lnot\left( \lnot A\right) \Longleftrightarrow A$ (llei de la doble
negació)

\item $A\wedge B\Longleftrightarrow B\wedge A$ (llei commutativa de $\wedge$)

\item $A\vee B\Longleftrightarrow B\vee A$ (llei commutativa de $\vee$)

\item $\left( A\wedge B\right) \wedge C\Longleftrightarrow A\wedge\left(
B\wedge C\right) $ (llei associativa de $\wedge$)

\item $\left( A\vee B\right) \vee C\Longleftrightarrow A\vee\left( B\vee
C\right) $ (llei associativa de $\vee$)

\item $A\vee\left( B\wedge C\right) \Longleftrightarrow\left( A\vee B\right)
\wedge\left( A\vee C\right) $ (llei distributiva de $\vee$ respecte de $\wedge$)

\item $A\wedge\left( B\vee C\right) \Longleftrightarrow\left( A\wedge
B\right) \vee\left( A\wedge C\right) $ (llei distributiva de $\wedge$
respecte de $\vee$)

\item $A\longrightarrow B\Longleftrightarrow\lnot A\vee B$

\item $A\longrightarrow B\Longleftrightarrow\lnot B\longrightarrow\lnot A$
(llei del contrarecíproc)

\item $A\longleftrightarrow B\Longleftrightarrow\left( A\longrightarrow
B\right) \wedge\left( B\longrightarrow A\right) $ (llei del bicondicional)

\item $\left( A\vee B\right) \longrightarrow C\Longleftrightarrow\left(
A\longrightarrow C\right) \wedge\left( B\longrightarrow C\right) $

\item $A\longrightarrow\left( B\wedge C\right) \Longleftrightarrow\left(
A\longrightarrow B\right) \wedge\left( A\longrightarrow C\right) $

\item $\lnot(A\wedge B)\Longleftrightarrow\lnot A\vee\lnot B$ (llei de De
Morgan)

\item $\lnot(A\vee B)\Longleftrightarrow\lnot A\wedge\lnot B$ (llei de De
Morgan)
\end{enumerate}

Algunes d'aquestes lleis tenen molt d'interès quan volem fer
demostracions. La llei de la doble negació a la pràctica es usa de
la següent manera: volem provar que $p$ és vertadera. Per a fer-ho,
suposem com hipòtesi que $\lnot p$ és vertadera, i deduïm
contradicció, i, per tant, $\lnot\left( \lnot p\right) $ és
vertadera. Ara bé, $\lnot\left( \lnot p\right) $ i $p$ són
equivalents i, per tant, $p$ és vertadera. Aquesta manera de demostrar
un enunciat matemàtic serà tractada amb més detall més
endavant.

\bigskip

Es diu \textbf{recíproc} de la proposició $p\longrightarrow q$ a la
proposició $q\longrightarrow p$, i \textbf{contrarecíproc} a la
proposició $\lnot q\longrightarrow\lnot p$. La llei del contrarecíproc també és la base de moltes demostracions. A la pràctica es
usa de la següent manera: volem provar que $p$ implica $q$. Per a
fer-ho, observem que és més senzill provar que $\lnot q$ implica $\lnot p$, i, per tant, és cert que $p$ implica $q$. Tractarem aquesta
forma de demostració més en detall més endavant.

\bigskip

Finalment, fem alguns comentaris sobre la llei del bicondicional. Quan s'ha
de demostrar una equivalencia $p\longleftrightarrow q$, la llei del
bicondicional permet fer-ho de la següent manera: primer, provem que $p$
implica $q$, i després, el recíproc, es a dir, que $q$ implica $p$.
Es habitual en matemàtiques dir-ho d'alguna d'aquestes maneres: (1) $p$
és condició necessària i suficient per $q$ o (2) si $p$, llavors
$q$, i recíprocament.

\bigskip

Per a completar aquesta secció provarem les lleis 10 i 12. Per demostrar
10 hem de comprovar que $\left( A\longleftrightarrow B\right)
\longleftrightarrow\left( A\longrightarrow B\right) \wedge\left(
B\longrightarrow A\right) $ és tautologia. En efecte, constriuïm la
matriu i s'obté tautologia

\begin{equation*}
\begin{tabular}{ccccccc}
$A$ & $B$ & $\alpha=A$ $\longrightarrow$ $B$ & $\beta=B\longrightarrow A$ & $\gamma=A\longleftrightarrow B$ & $\alpha\wedge\beta$ & $\gamma
\longleftrightarrow\left( \alpha\wedge\beta\right) $ \\ \hline
V & V & V & V & V & V & V \\
V & F & F & V & F & F & V \\
F & V & V & F & F & F & V \\
F & F & V & V & V & V & V\end{tabular}
\text{.}
\end{equation*}

Per demostrar 11 hem de comprovar que $\lnot(A\wedge B)\longleftrightarrow
\left( \lnot A\vee\lnot B\right) $ és tautologia. En efecte, constriuïm la matriu i s'obté tautologia

\begin{equation*}
\begin{tabular}{llllllll}
$A$ & $B$ & $A\wedge B$ & $\lnot A$ & $\lnot B$ & $\alpha=\lnot\left(
A\wedge B\right) \,$ & $\beta=\lnot A\vee\lnot B$ & $\alpha\longleftrightarrow\beta $ \\ \hline
\multicolumn{1}{c}{V} & \multicolumn{1}{c|}{V} & \multicolumn{1}{c}{V} &
\multicolumn{1}{c}{F} & \multicolumn{1}{c}{F} & \multicolumn{1}{c}{F} &
\multicolumn{1}{c}{F} & \multicolumn{1}{c}{V} \\
\multicolumn{1}{c}{V} & \multicolumn{1}{c|}{F} & \multicolumn{1}{c}{F} &
\multicolumn{1}{c}{F} & \multicolumn{1}{c}{V} & \multicolumn{1}{c}{V} &
\multicolumn{1}{c}{V} & \multicolumn{1}{c}{V} \\
\multicolumn{1}{c}{F} & \multicolumn{1}{c|}{V} & \multicolumn{1}{c}{F} &
\multicolumn{1}{c}{V} & \multicolumn{1}{c}{F} & \multicolumn{1}{c}{V} &
\multicolumn{1}{c}{V} & \multicolumn{1}{c}{V} \\
\multicolumn{1}{c}{F} & \multicolumn{1}{c|}{F} & \multicolumn{1}{c}{F} &
\multicolumn{1}{c}{V} & \multicolumn{1}{c}{V} & \multicolumn{1}{c}{V} &
\multicolumn{1}{c}{V} & \multicolumn{1}{c}{V}\end{tabular}
.
\end{equation*}

Les altres és demostran de la mateixa manera.

\begin{exemple}
Escriu fent ús de les connectives les proposicions següents: (1) un nombre
natural és parell si i només si el seu quadrat és parell; (2) és necessari
que $x$ sigui un nombre real no negatiu perquè $\sqrt{x}$ sigui real; (3) és
suficient que $n=3$ perquè $n^{2}-5n+6=0$; (4) $n^{2}-5n+6=0$ precisament si
$n=2$ o $n=3$; (5) $\left\vert a\right\vert=\left\vert b\right\vert$ si i
només si $a=b$ o $a=-b$.
\end{exemple}

\begin{solucio}
(1) Si $P(n)$ vol dir \enquote{$n$ és parell}, l'enunciat és
\[
P(n)\longleftrightarrow P(n^{2}).
\]
(2) Si $R(x)$ vol dir \enquote{$\sqrt{x}$ és real}, l'enunciat és
\[
R(x)\longrightarrow (x\in\mathbb{R}\wedge x\geq0).
\]
(3) L'enunciat és
\[
n=3\longrightarrow n^{2}-5n+6=0.
\]
(4) L'enunciat és
\[
n^{2}-5n+6=0\longleftrightarrow(n=2\vee n=3).
\]
(5) L'enunciat és
\[
\left\vert a\right\vert=\left\vert b\right\vert\longleftrightarrow(a=b\vee a=-b).
\]
\end{solucio}

\section{Inferència lògica}

Molts cops veieu com en resoldre exercicis a matemàtiques se us demana
que raoneu o justifiqueu la resposta. Ara volem tractar aquest assumpte. Què es vol dir que raoneu la vostra resposta? És clar, heu de fer un
raonament vàlid que justifiqui allò que se us demana demostrar. Això té dues parts, una és conduir correctament el vostre raonament
des del punt vista lògic, i l'altra, és comprendre correctament els
conceptes que estan presents en el raonament des del punt de vista matemàtic. Aquí tractarem el primer punt perquè el segon depèn dels
coneixements que cadascú té de les matemàtiques.

\bigskip

Si suposem que la proposició condicional \enquote{$p\rightarrow q$} és
vertadera, llavors sabem (recordeu la taula de veritat d'aquesta connectiva)
que si \enquote{$p$} és vertadera, necessàriament ho serà \enquote{$q$}. Ara bé, només pel fet que \enquote{$p\rightarrow q$} sigui vertadera no es té
que \enquote{$p$} i \enquote{$q$} siguin vertaderes, perquè podrien ser totes dues
falses, o \enquote{$p$} falsa i \enquote{$q$} vertadera. Per tant, si \enquote{$p\rightarrow q$} i \enquote{$p$} són vertaderes, llavors sí que \enquote{$q$} ha de ser vertadera. En
aquest últim cas es diu que \enquote{$q$} és \textbf{conseqüència lògica} de \enquote{$p\rightarrow q$} i \enquote{$p$} o també que \enquote{$q$} s'\textbf{infereix lògicament} de \enquote{$p\rightarrow q$} i \enquote{$p$}. És habitual
representar aquesta inferència lògica seguint el següent format:\begin{equation*}
\begin{tabular}{l}
$p\longrightarrow q$ \\
$p$ \\ \hline
$q$\end{tabular}
\ \text{,}
\end{equation*}
on la linea horitzontal separa les proposicions que es diuen \textbf{premisses}, de la proposició que es diu \textbf{conclusió}.

Per exemple: Si $x\in A$, llavors $\left\vert x\right\vert \leq1$; sabem que
$x\in A\cap B$. En conseqüència, $\left\vert x\right\vert \leq1$. La
regla d'inferència anterior ens permet assegurar que aquest argument
és vàlid perquè sabem que $x\in A\cap B$ implica $x\in A$.

\bigskip

Observem doncs, que la conseqüència lògica és una relació
entre les premisses i la conclusió. De fet, pel raonament que hem fet
abans aquesta inferència pot també escriure's com \enquote{$\left(
p\rightarrow q\right) \wedge p$} $\Longrightarrow$ \enquote{$q$}, és a dir, com
una implicació lògica. És clar que també podríem
reescriure aquesta implicació lògica en termes de formes
proposicionals com $(A\longrightarrow B)\wedge A$ $\Longrightarrow B$, on $A$
i $B$ són formes, i també com a inferència lògica seguint el
format anterior:\begin{equation*}
\begin{tabular}{l}
$A\longrightarrow B$ \\
$A$ \\ \hline
$B$\end{tabular}
\text{.}
\end{equation*}

En general, des del punt de vista lògic un \textbf{raonament} és un
condicional que té com antecedent la conjunció de proposicions $P_{1},...,P_{n}$, anomenades premisses, i com a conseqüent una proposició $C$, anomenada conclusió:\begin{equation*}
\left( P_{1}\wedge\cdots\wedge P_{n}\right) \longrightarrow C\text{.}
\end{equation*}
Es diu que el raonament és \textbf{vàlid} si la conclusió necessàriament es deriva de les premisses, o sigui, si $\left(
P_{1}\wedge\cdots\wedge P_{n}\right) \longrightarrow C$ és tautologia, o
equivalentment, si $\left( P_{1}\wedge\cdots\wedge P_{n}\right)
\Longrightarrow C$. Pensant en la noció d'implicació lògica,
podem dir que un raonament és vàlid si no podem assignar valors de
veritat a les proposicions que s'utilitzen en el raonament de manera que les
premisses siguin vertaderes i la conclusió sigui falsa.

\bigskip

És important no oblidar aquí que la lògica s'ocupa només
d'analitzar la validesa dels raonaments, o sigui de la sintaxi del
llenguatge, no ens pot dir res sobre si la informació continguda en una
hipòtesi és vertadera o falsa, que formaria part de la interpretació o semàntica del llenguatge. Per tant, els termes vàlid i no vàlid es refereixen a l'estructura del raonament, no a la veritat o
falsedat de les proposicions què depèn dels nostres coneixements de
matemàtiques.

\bigskip

Tenint en compte l'última manera d'escriure el nostre argument, podríem intentar demostrar que és vàlid, mostrant que $\left(
P_{1}\wedge\cdots\wedge P_{n}\right) \longrightarrow C$ és tautologia,
utilitzant una taula de veritat. Aquest mètode realment funcionaria, però no seria gens eficaç. En primer lloc, atès que si hi ha $m$
proposicions implicades, la taula de veritat hauria de tenir $2^{m}$ files,
la qual cosa seria molt feixuga si $m$ és gran. En segon lloc, l'ús
d'una taula de veritat no proporciona cap visió intuïtiva de per què l'argument és vàlid.

Per a aquests motius, en lloc d'utilitzar taules de veritat, intentarem
justificar la validesa dels raonaments fent ús de les implicacions lògiques que hem donat. Si volem demostrar una implicació lògica
complicada, és recomanable fer-ho descomponent-la en una
col·lecció d'implicacions més senzilles, preses
d'una en una. Si ja es coneixen les implicacions més simples, podrien
ser blocs per a la implicació més complicada. Algunes de les
implicacions senzilles que fem servir, conegudes com a \textbf{regles d'inferència} \label{sec:lleis-inferencia} de la lògica proposicional, es detallen a continuació.

\bigskip

\begin{center}
\resizebox{\textwidth}{!}{\begin{tabular}{lllll}
\begin{tabular}{l}
$A\longrightarrow B$ \\
$A$ \\ \hline
$B$\end{tabular}
& MP: Modus Ponens &  &
\begin{tabular}{l}
$A\longrightarrow B$ \\
$\lnot B$ \\ \hline
$\lnot A$\end{tabular}
& MT: Modus Tollens \\
&  &  &  &  \\
\begin{tabular}{l}
$A$ \\ \hline\hline
$\lnot\lnot A$\end{tabular}
& DN: Doble Negació &  &
\begin{tabular}{l}
$A$ \\ \hline\hline
$A$\end{tabular}
& R: Repetició \\
&  &  &  &  \\
\begin{tabular}{l}
$A\wedge B$ \\ \hline
$A$\end{tabular}
& EC: Eliminació Conjunctor &  &
\begin{tabular}{l}
$A\wedge B$ \\ \hline
$B$\end{tabular}
& EC: Eliminació Conjunctor \\
&  &  &  &  \\
\begin{tabular}{l}
$A$ \\
$B$ \\ \hline
$A\wedge B$\end{tabular}
& IC: Introducció conjunctor &  &
\begin{tabular}{l}
$A$ \\ \hline
$A\vee B$\end{tabular}
& IC: Introducció disjunctor \\
&  &  &  &  \\
\begin{tabular}{l}
$B$ \\ \hline
$A\vee B$\end{tabular}
& ID: Introducció Disjunctor &  &
\begin{tabular}{l}
$A\vee B$ \\
$\lnot A$ \\ \hline
$B$\end{tabular}
& ED: Eliminació Disjunctor \\
&  &  &  &  \\
\begin{tabular}{l}
$A\vee B$ \\
$\lnot B$ \\ \hline
$A$\end{tabular}
& ED: Eliminació Disjunctor &  &
\begin{tabular}{l}
$A\longleftrightarrow B$ \\ \hline
$A\longrightarrow B$\end{tabular}
& EB: Eliminació Bicondicional \\
&  &  &  &  \\
\begin{tabular}{l}
$A\longleftrightarrow B$ \\ \hline
$B\longrightarrow A$\end{tabular}
& EB: Eliminació Bicondicional &  &
\begin{tabular}{l}
$A\longrightarrow B$ \\
$B\longrightarrow A$ \\ \hline
$A\longleftrightarrow B$\end{tabular}
& IB: Introducció Bicondicional \\
&  &  &  &  \\
\begin{tabular}{l}
$A\longrightarrow B$ \\
$B\longrightarrow C$ \\ \hline
$A\longrightarrow C$\end{tabular}
& TC: Transitivitat Condicional &  &
\begin{tabular}{l}
$A\longrightarrow B$ \\
$C\longrightarrow D$ \\
$A\vee C$ \\ \hline
$B\vee D$\end{tabular}
& DC: Dilema constructiu \\
&  &  &  &  \\
\begin{tabular}{l}
$\lnot\left( A\vee B\right) $ \\ \hline\hline
$\lnot A\wedge\lnot B$\end{tabular}
& DM: De Morgan &  &
\begin{tabular}{l}
$\lnot\left( A\wedge B\right) $ \\ \hline\hline
$\lnot A\vee\lnot B$\end{tabular}
& DM: De Morgan\end{tabular}}
\end{center}

\bigskip

En aquesta taula hem utilitzat el conveni de representar per una línia
horitzontal doble dues regles d'inferència com s'indica a continuació:\begin{equation*}
\begin{tabular}{l}
$A$ \\ \hline\hline
$B$\end{tabular}
\ \qquad\text{està en lloc de\qquad\ }\begin{tabular}{l}
$A$ \\ \hline
$B$\end{tabular}
\ \quad\text{ i \quad}\begin{tabular}{l}
$B$ \\ \hline
$A$\end{tabular}
\
\end{equation*}
on $A$ i $B$ són formes proposicionals. No sols hem tingut present les
implicacions lògiques de la secció \ref{sec:implicacio-logica}, també les equivalències lògiques de la secció \ref{sec:equivalencia-logica} com, per exemple, les lleis
de De Morgan.

\begin{exemple}
Volem determinar si el següent raonament és vàlid:
\enquote{Si en Pau va rebre l'e-mail, llavors va agafar l'avió
i serà aquí al migdia. Pau no va agafar l'avió. Per tant, Pau
no va rebre l'e-mail}.
\end{exemple}

\begin{solucio}
Primer procedim a identificar les proposicions simples:

\begin{description}
\item[$p=$] \enquote{Pau va rebre l}e-mail'

\item[$q=$] \enquote{Pau va agafar l}avió'

\item[$r=$] \enquote{Pau serà aquí al migdia}
\end{description}

Després, expressem formalment el raonament:\begin{equation*}
\begin{tabular}{ll}
$P_{1}:$ & $p\longrightarrow\left( q\wedge r\right) $ \\
$P_{2}:$ & $\lnot q$ \\ \hline
$C:$ & $\lnot p$\end{tabular}
\end{equation*}

Per provar la validesa d'aquest argument utilitzarem les regles d'inferència:\begin{equation*}
\begin{tabular}{lll}
1. & $p\longrightarrow\left( q\wedge r\right) $ & $P_{1}$ \\
2. & $\lnot q$ & $P_{2}$ \\
3. & $\lnot q\vee\lnot r$ & ID 2 \\
4. & $\lnot\left( q\wedge r\right) $ & DM 3 \\ \hline
5. & $\lnot p$ & MT (1,4)\end{tabular}
\
\end{equation*}
Observem que la fila $3$ s'ha deduït de la fila 2 mitjançant la
regla ID; la fila 4 s'ha deduït de la fila 3 fent ús de la llei de
De Morgan i, finalment, la fila 5 que és la conclusió s'ha obtingut
de les files 1 i 4 per la regla MT.
\end{solucio}

\begin{exemple}
Volem determinar si el següent raonament és vàlid: "Si robes un banc, vas a
la presó. Si vas a la presó, no et diverteixes. Si tens vacances, et
diverteixes. Robes un banc o tens vacances. Per tant, vas a la presó o et
diverteixes".
\end{exemple}

\begin{solucio}
Les proposicions simples d'aquest raonament són:

\begin{description}
\item[$p=$] \enquote{Robes un banc}

\item[$q=$] \enquote{Vas a la presó}

\item[$r=$] \enquote{Et diverteixes}

\item[$s=$] \enquote{Tens vacances}
\end{description}

Expressem el raonament simbòlicament:
\begin{equation*}
\begin{tabular}{ll}
$P_{1}:$ & $p\longrightarrow q$ \\
$P_{2}:$ & $q\longrightarrow\lnot r$ \\
$P_{3}:$ & $s\longrightarrow r$ \\
$P_{4}:$ & $p\vee s$ \\ \hline
$C:$ & $q\vee r$\end{tabular}
\end{equation*}

El raonament és vàlid. En efecte, de $p\longrightarrow q$, de
$s\longrightarrow r$ i de $p\vee s$ s'obté $q\vee r$ per dilema constructiu:
\begin{equation*}
\begin{tabular}{lll}
1. & $p\longrightarrow q$ & $P_{1}$ \\
2. & $q\longrightarrow\lnot r$ & $P_{2}$ \\
3. & $s\longrightarrow r$ & $P_{3}$ \\
4. & $p\vee s$ & $P_{4}$ \\ \hline
5. & $q\vee r$ & DC (1,3,4)\end{tabular}
\end{equation*}
La premissa $P_{2}$ no s'utilitza en aquesta deducció, però això no invalida
l'argument. Una premissa addicional pot ser redundant; el que no pot passar,
en un raonament vàlid, és que totes les premisses siguin vertaderes i la
conclusió falsa.
\end{solucio}

\bigskip

En tot raonament se'ns presenten dues qüestions molt importants: Una
és la seva validesa, i l'altra, és la seva deduïbilitat. La
primera qüestió fa referència a les taules de veritat, i la
segona, a les regles d'inferència. Quina relació hi ha entre
aquestes dues nocions? Tot i que no és absolutament obvi ni fàcil de
demostrar, resulta molt remarcable que les dues nocions, una de naturalesa
semàntica i l'altre sintàctica, sempre donen el mateix resultat,
és a dir, un raonament és vàlid si i només si és deduïble. Per tant, si volem demostrar que un argument donat és vàlid, n'hi haurà prou amb demostrar que és deduïble i viceversa.
L'equivalència d'aquests dos enfocaments és un resultat important en
la lògica. El fet que la validesa implica la deduïbilitat es coneix
com a \textbf{teorema de completesa} de la lògica proposicional, i el
fet que la deduïbilitat implica la validesa es coneix com a \textbf{teorema de la correcció} de la lògica proposicional. A més,
aquesta lògica és \textbf{decidible}, la qual cosa vol dir que hi ha
un procediment finit, per exemple les taules de veritat, que permet determinar
si qualsevol raonament expressat en aquest llenguatge és vàlid o no. Aquests
resultats i les seves demostracions poden trobar-se
en qualsevol text de lògica proposicional.

\bigskip

Per les consideracions de l'apartat anterior observem que per provar que un
raonament és vàlid, simplement hem de trobar una deducció, que
sovint és una tasca molt més agradable que mostrar directament la
seva validesa. En canvi, per provar que un raonament no és vàlid,
les deduccions no són de gran ajuda, perquè hauríem de
demostrar que no és possible trobar cap deducció, i això no
podrem assegurar-ho mai, sempre podria haver-hi una deducció que funcionés. En aquests casos és millor fer servir la definició de
validesa directament i trobar valors de veritat per als quals les
proposicions que són premisses del raonament siguin vertaderes i la
conclusió falsa.

\begin{exemple}
Volem determinar la validesa del següent raonament: \enquote{Si
el crim va ocórrer després de les 4, llavors en Pep no va poder
haver-lo comès. Si el crim va ocórrer a les 4 o abans, llavors en
Carles no va poder haver-lo comès. El crim involucra a dues persones, si
en Carles no el va cometre. Per tant, el crim involucra a dues
persones}.
\end{exemple}

\begin{solucio}
Primer procedim a identificar les proposicions simples:

\begin{description}
\item[$p=$] \enquote{El crim va ocórrer després de les 4}

\item[$q=$] \enquote{Pep podia haver comès el crim}

\item[$r=$] \enquote{Carles podia haver comès el crim}

\item[$s=$] \enquote{El crim involucra a dues persones}
\end{description}

Després expressem el raonament simbòlicament:\begin{equation*}
\begin{tabular}{ll}
$P_{1}:$ & $p\longrightarrow\lnot q$ \\
$P_{2}:$ & $\lnot p\longrightarrow\lnot r$ \\
$P_{3}:$ & $\lnot r\longrightarrow s$ \\ \hline
$C:$ & $s$\end{tabular}
\
\end{equation*}

Intentem provar la validesa d'aquest argument utilitzant les regles d'inferència:\begin{equation*}
\begin{tabular}{lll}
1. & $p\longrightarrow\lnot q$ & $P_{1}$ \\
2. & $\lnot p\longrightarrow\lnot r$ & $P_{2}$ \\
3. & $\lnot r\longrightarrow s$ & $P_{3}$ \\
4. & $\lnot p\longrightarrow s$ & TC (2,3) \\
5. & $p\vee\lnot p$ & Tautologia \\ \hline
6. & $\lnot q\vee s$ & DC (1,4,5)\end{tabular}
\
\end{equation*}
No hem deduït $s$ sinó $\lnot q\vee s$. Això fa pensar que si $q $ és falsa i $s$ també, hi ha una interpretació (una fila de
la taula de veritat) que farà les premisses vertaderes i la conclusió, falsa; es pot comprovar que això passa si prenem $p$ i $r$ ambdues
vertaderes. Per tant, aquest raonament no és vàlid. A continuació
mostrem la interpretació esmentada:\begin{equation*}
\begin{tabular}{l|l|l|l||l|l|l|l|l||l|l|l|l||l}
$p$ & $\longrightarrow$ & $\lnot$ & $q$ & $\lnot$ & $p$ & $\longrightarrow$
& $\lnot$ & $r$ & $\lnot$ & $r$ & $\longrightarrow$ & $s $ & $s$ \\ \hline
V & V & V & F & F & V & V & F & V & F & V & V & F & F\end{tabular}
\
\end{equation*}

Volem destacar el fet què construir una taula de veritat per veure que
hi ha una valoració que prova el que hem dit, portaria molta feina perquè la taula tindria 16 files.
\end{solucio}

\bigskip

Hi ha raonaments que quan intentem provar la seva validesa deduïm una
contradicció, és a dir, deduïm una proposició i la seva
negació. Llavors es diu que el conjunt de premisses d'aquest raonament
és \textbf{inconsistent}. Quan les premisses d'un raonament no són
inconsistents es diu que les premisses són \textbf{consistents}. No és que hi hagi res lògicament erroni amb premisses inconsistents,
senzillament que no serveixen per res, perquè aleshores podem deduir
qualsevol proposició. Volem destacar també que, des del punt de
vista semàntic, un raonament té un conjunt de premisses inconsistent
quan no hi ha cap interpretació (cap fila de la taula de veritat) que
faci a totes les premisses vertaderes.

\begin{exemple}
Volem determinar si les premisses del següent raonament és o no
consistent: "Miquel no toca la piano o Maria toca la guitarra. Si Carla no
toca el violí, Maria no toca la guitarra. Miquel toca el piano i Carla
no toca el violí. Per tant, Jaume toca l'acordió".
\end{exemple}

\begin{solucio}
Procedim a formalitzar el raonament, identificant les seves proposicions
simples:

\begin{description}
\item[$p=$] \enquote{Miquel toca la piano}

\item[$q=$] \enquote{Maria toca la guitarra}

\item[$r=$] \enquote{Carla toca el violí}

\item[$s=$] \enquote{Jaume toca l}acordió'
\end{description}

Llavors el raonament formalitzat és\begin{equation*}
\begin{tabular}{ll}
$P_{1}:$ & $\lnot p\vee q$ \\
$P_{2}:$ & $\lnot r\longrightarrow\lnot q$ \\
$P_{3}:$ & $p\wedge\lnot r$ \\ \hline
$C:$ & $s$\end{tabular}
\
\end{equation*}

Per un costat, utilitzant les regles d'inferència, s'obté:\begin{equation*}
\begin{tabular}{lll}
1. & $\lnot p\vee q$ & $P_{1}$ \\
2. & $\lnot r\longrightarrow\lnot q$ & $P_{2}$ \\
3. & $p\wedge\lnot r$ & $P_{3}$ \\
4. & $p$ & EC 3 \\
5. & $p\vee s$ & ID 4 \\
6. & $\lnot r$ & EC 3 \\
7. & $\lnot q$ & MP (2,6) \\
8. & $\lnot p$ & ED (1,7) \\ \hline
9. & $s$ & ED (5,8)\end{tabular}
\
\end{equation*}
Per un altre, es té també\begin{equation*}
\begin{tabular}{lll}
1. & $\lnot p\vee q$ & Hipòtesis \\
2. & $\lnot r\longrightarrow\lnot q$ & Hipòtesis \\
3. & $p\wedge\lnot r$ & Hipòtesis \\
4. & $p$ & EC 3 \\
5. & $p\vee\lnot s$ & ID 4 \\
6. & $\lnot r$ & EC 3 \\
7. & $\lnot q$ & MP (2,6) \\
8. & $\lnot p$ & ED (1,7) \\ \hline
9. & $\lnot s$ & ED (5,8)\end{tabular}
\
\end{equation*}
Observem que hem deduït $s$ en el primer cas, i $\lnot s$, en el segon.
Com a conseqüència, podem deduïr $s\wedge\lnot s$, que és una
contradicció. Per tant, el conjunt de premisses és inconsistent.

També podem provar la inconsistència de les premisses comprovant si
és possible que hi hagi una interpretació que faci totes les
premisses vertaderes:
\begin{equation*}
\begin{tabular}{|l|l|l|l|l|l|l|l|l|l|l|l|l|}
$\lnot$ & $p$ & $\vee$ & $q$ & $\lnot$ & $r$ & $\longrightarrow$ & $\lnot$ &
$q$ & $p$ & $\wedge$ & $\lnot$ & $r$ \\ \hline
&  & V &  &  &  & V ! &  &  &  & V &  &  \\ \cline{7-7}
F & V &  & V & V & F & F ! & F & V & V &  & V & F \\ \cline{7-7}
\end{tabular}
\
\end{equation*}
Observem que no és possible excepte que acceptem que una proposició
sigui vertadera i falsa alhora, cosa que no pot ser.
\end{solucio}

\bigskip

Per acabar aquesta secció, volem tractar els raonaments que a vegades es
consideren com a vàlids i no ho són, sovint coneguts com a \textbf{fal·làcies}. En primer lloc tractarem els que tenen a
veure amb un mal ús de les regles d'inferència. Són arguments
que poden semblar vàlids a primera vista, però que són
fal·làcies. Considerem el següent raonament: \enquote{Si
l'Albert menja un bon dinar, beurà una cervesa. L'Albert va beure una
cervesa i, per tant, va menjar un bon dinar}. Podríem pensar que aquest
raonament és vàlid per MP, però no és així. En efecte,
si formalitzem l'argument es té:

\begin{description}
\item[$p=$] \enquote{Albert menja un bon dinar}

\item[$q=$] \enquote{Albert beurà una cervesa}
\end{description}

i, aleshores,\begin{equation*}
\begin{tabular}{ll}
$P_{1}:$ & $p\longrightarrow q$ \\
$P_{2}:$ & $q$ \\ \hline
$C:$ & $p$\end{tabular}
\end{equation*}
És evident que no poden aplicar MP i, a part, no és vàlid, perquè pot haver begut una cervesa sense haver tingut un bon sopar. De fet,
si fem la interpretació següent\begin{equation*}
\begin{tabular}{lll||l||l}
$p$ & $\longrightarrow$ & $q$ & $q$ & $p$ \\ \hline
F & V & V & V & F\end{tabular}
\end{equation*}
comprovem el que hem dit.

\bigskip

Considerem ara el següent raonament: "Si en Joan va en bicicleta a
primera hora del matí, aleshores menja un bon esmorzar. En Joan no va
en bicicleta. Per tant, en Joan no menja un bon esmorzar". Podríem
pensar que aquest raonament és vàlid per MT, però no és així. En efecte, si formalitzem l'argument es té:

\begin{description}
\item[$p=$] \enquote{Joan va en bicicleta a primera hora del matí}

\item[$q=$] \enquote{Joan menja un bon esmorzar}
\end{description}

i, aleshores,\begin{equation*}
\begin{tabular}{ll}
$P_{1}:$ & $p\longrightarrow q$ \\
$P_{2}:$ & $\lnot p$ \\ \hline
$C:$ & $\lnot q$\end{tabular}
\end{equation*}
És evident que no poden aplicar MT i, a part, no és vàlid, perquè pot haver menjat un bon esmorzar sense haver anat en bicicleta. De
fet, si fem la interpretació següent\begin{equation*}
\begin{tabular}{lll||l|l||ll}
$p$ & $\longrightarrow$ & $q$ & $\lnot$ & $p$ & $\lnot$ & $q$ \\ \hline
F & V & V & V & F & F & V\end{tabular}
\end{equation*}
comprovem el que hem dit.

\bigskip

En segon lloc, hi ha raonaments que no són vàlids perquè es fan
suposicions que no estan justificades per les premisses. Per exemple,
considerem el següent raonament: "Si en Dídac té febre,
esternuda molt. Per tant, Dídac esternuda molt". És clar que podríem aplicar MP per concloure que en Dídac esternuda molt, però
no sabem per les hipòtesis que en Dídac té febre.

Aquests exemples de fal·làcies poden semblar molt
senzills, però quan el raonament és més extens i complex, i no
s'escriu sinó que es parla, a vegades aquests errors lògics passen
desapercebuts. Cal doncs tenir cura de no cometre'ls.

\section{Quantificadors}

En la secció \enquote{Variables i constants} hem escrit
\begin{equation*}
x+y=y+x\text{,}
\end{equation*}
reconeixent en aquesta expressió la llei commutativa de l'aritmètica. Si assumim que les variables designen nombres reals, aquest enunciat
afirma que tots els nombres reals compleixen aquesta propietat i, en
general, tots els enunciats d'aquest tipus que afirmen que objectes
arbitraris d'una classe compleixen una determinada propietat, s'anomenen
\textbf{proposicions universals}.

A vegades també ens trobem enunciats com\begin{equation*}
x^{2}-3x+2=0
\end{equation*}
reconeixent en aquesta expressió que hi ha dos nombres que compleixen
aquesta equació. En aquests casos diem que són \textbf{proposicions
existencials} perquè afirmen l'existència d'objectes (en aquest cas,
els nombres 1 i 2) que compleixen una certa propietat.

\bigskip

Fins ara, mitjançant els símbols lògics $\wedge,\vee
,\lnot,\longrightarrow$ i $\longleftrightarrow$, hem vist com podem
formalitzar molts enunciats de les matemàtiques i com podem conèixer
que dos enunciats diferents tenen el mateix significat si són
equivalents. També hem vist com la formalització ens pot ajudar a
comprendre l'estructura lògica dels raonaments, i com els raonaments són vàlids amb l'ajut d'implicacions lògiques o regles d'inferència. Però amb tot això no n'hi ha prou per captar la totalitat
del significat de molts enunciats de les matemàtiques. A l'inici
d'aquesta secció hem vist com hi ha uns enunciats molt comuns a matemàtiques que cal estudiar amb més detall.

\bigskip

En efecte, imaginem que estem tractant un conjunt infinit $X$ de nombres
enters. Considerem ara l'enunciat \enquote{Tots els elements de $X$ són senars}.
Si volem formalitzar aquest enunciat, hauríem d'escriure\begin{equation*}
P(x_{1})\wedge P(x_{2})\wedge P(x_{3})\wedge\cdots\text{,}
\end{equation*}
on $P(x)$ és el predicat \enquote{$x$ és senar}. I si volem formalitzar
l'enunciat \enquote{Hi ha almenys un element de $X$ que és senar}, hauríem
d'escriure\begin{equation*}
P(x_{1})\vee P(x_{2})\vee P(x_{3})\vee\cdots\text{.}
\end{equation*}
El problema d'aquestes dues expressions és que no s'acaben mai. Per
superar-lo introduirem dos nous símbols $\exists$ i $\forall$. El símbol $\exists$, anomenat \textbf{quantificador existencial}, està
en lloc de frases com \enquote{existeix} o \enquote{hi ha un}. D'aquesta manera l'enunciat
\enquote{Hi ha elements de $X$ que són senars} podem escriure'l així\begin{equation*}
\exists x\in X,P(x)\text{.}
\end{equation*}

El símbol $\forall$, anomenat \textbf{quantificador universal}, està
en lloc de frases com \enquote{per a tots} o \enquote{per a cadascun}. D'aquesta manera
l'enunciat \enquote{Tots els elements de $X$ són senars} podem escriure'l així\begin{equation*}
\forall x\in X,P(x)\text{.}
\end{equation*}

\bigskip

Els exemples que hem donat a l'inici d'aquesta secció podem ara
expressar-los d'aquesta manera:
\begin{equation*}
\left( \forall x,y\right) \left( x,y\in\mathbb{R}\longrightarrow
x+y=y+x\right)
\end{equation*}
que constitueix una proposició universal, i l'altre\begin{equation*}
\left( \exists x\right) \left( x\in\mathbb{R}\wedge x^{2}-3x+2=0\right)
\end{equation*}
que és una proposició existèncial.

\begin{observacio}
En teoria de conjunts la proposició $\exists x\in X,P(x)$ podem
escriure-la d'aquesta manera:\begin{equation*}
\left\{ x\in X:P(x)\right\} \neq\emptyset\text{,}
\end{equation*}
indicant que el conjunt dels elements que compleixen la propietat $P$ és
no buit; i la proposició $\forall x\in X,P(x)$, com\begin{equation*}
\left\{ x\in X:P(x)\right\} =X\text{,}
\end{equation*}
indicant que el conjunt dels elements que compleixen la propietat $P$ és
tot el conjunt $X$.
\end{observacio}

\bigskip

L'enunciat \enquote{el quadrat d}un nombre real més petit o igual que 4' és
un predicat. Si $x$ és un nombre real, aleshores \enquote{$x^{2}\leq4$} és
el predicat abreujat. L'enunciat \enquote{$x^{2}\leq4$} no és una proposició
perquè no podem afirmar que sigui vertadera o falsa llevat que assignem
a la varible lliure $x$ un valor real. Si representem per $P(x)$ aquest
predicat, llavors $P(1)$ és una proposició vertadera, en canvi, $P(3) $ és falsa.

Hi ha una altra manera de fer que un predicat es converteixi en una proposició. Es fent que les seves variables lliures quedin lligades per
quantificadors. Considerem que $P(x)$ el predicat anterior, aleshores \enquote{$\forall x\in\mathbb{R}$, $P(x)$} és una proposició falsa, i en
canvi, \enquote{$\exists x\in\mathbb{R}$, $P(x)$} és vertadera.

\bigskip

Podem construir enunciats amb més d'un quantificador. Per exemple,
considerem el predicat $Q(m,n)=$ \enquote{$m<n$}, on $m,n\in\mathbb{Z}$. Llavors, la
proposició \enquote{$\forall m\in\mathbb{Z}~\exists n\in\mathbb{Z},m<n$},
escrita en català com \enquote{per a tot nombre enter $m$ existeix algun altre
enter $n$ tal que $m<n$}, és vertadera. En efecte, donat qualsevol enter
$m$, existeix $n=m+1$, també és enter i compleix la condició $m<n.$

\bigskip

L'ordre dels quantificadors és molt important perquè per exemple,
canviant l'ordre de l'enunciat anterior, o sigui \enquote{$\exists m\in\mathbb{Z}~\forall n\in\mathbb{Z},m<n$} es té una proposició falsa. És
important també observar el paper que fan les variables lliures en una
proposició que també en té de lligades. Amb relació al
predicat anterior, considerem la proposició \enquote{$\exists m\in\mathbb{Z},$ $m<n$}, què és vertadera perquè només cal prendre $m=n-1$. Si
ara posem $x$ en lloc de $m$ en la proposició, es té \enquote{$\exists x\in\mathbb{Z},$ $x<n$} que té el mateix significat que l'anterior; de fet,
els valors de $m$ o $x$ depenen sempre del valor que assignem a la variable
lliure $n$, i, per tant, si aquesta expressió formés part d'una
expressió més llarga, aquesta última no canviaria en el seu
significat. En canvi, si substituïm la variable lliure $n$ per una
altra $q$, la proposició \enquote{$\exists m\in\mathbb{Z},$ $m<q$} pot canviar
el significat d'una expressió més llarga de la qual en forma part,
doncs ara $m$ depèn de $q$ i no de $n$.

\bigskip

Ara volem veure la negació de proposicions amb quantificadors. Per
exemple, considerem el predicat $Q=$ \enquote{tenir els ulls blaus}. Llavors, si
suposem que la variable $x$ té per domini el conjunt de totes les
persones, l'enunciat \enquote{Totes les persones tenen els ulls blaus} s'escriu així: $\left( \forall x\right) Q(x)$. La negació d'aquest enunciat
és \enquote{No totes les persones tenen els ulls blaus} que s'escriu com $\lnot\left( \forall x\right) Q(x)$. És clar que aquest últim
enunciat també el podem escriure com \enquote{Hi ha alguna persona que no te els
ulls blaus}, que simbòlicament ho escrivim així: $\left( \exists
x\right) \lnot Q\left( x\right).$

Considerem ara el cas d'una proposició existencial. Considerem
l'enunciat \enquote{Existeix una solució real en l}equació $x^{3}+x=0$'. Si
suposem que el domini de la variable $x$ són els nombres reals,
aleshores escrivim aquesta proposició així: $\left( \exists
x\right) \left( x^{3}+x=0\right).$ La negació d'aquesta proposició
és \enquote{Cap nombre real és solució de l}equació $x^{3}+x=0$',
que s'escriu així: $\left( \forall x\right) \left( x^{3}+x\neq0\right).$

\bigskip

Podem resumir els dos casos anteriors d'aquesta manera: Si $P(x)$ és un
predicat i el domini de la variable $x$ és el conjunt $U$, aleshores es
compleixen les següents equivalències:

\begin{itemize}
\item $\lnot\left( \forall x\in U,P(x)\right) \Longleftrightarrow\exists
x\in U,\lnot P\left( x\right) $

\item $\lnot\left( \exists x\in U,P\left( x\right) \right)
\Longleftrightarrow\forall x\in U,\lnot P\left( x\right) $
\end{itemize}

\bigskip

A diferència de les equivalències discutides a la secció \ref{sec:equivalencia-logica}, no podem utilitzar taules de veritat per verificar aquestes equivalències, tot i que són certes, basant-se en els significats dels
quantificadors. Podem utilitzar les equivalències anteriors per negar
proposicions amb més d'un quantificador. Per exemple, suposem que $f$
és una funció real de variable i considerem l'enunciat \enquote{Per a cada
nombre real $x$, existeix un nombre real $y$ tal que $f(x)=y$}; simbòlicament s'escriu com $\left( \forall x\right) \left( \exists y\right)
\left( f(x)=y\right).$ La seva negació és\begin{align*}
\lnot\left( \left( \forall x\right) \left( \exists y\right) \left(
f(x)=y\right) \right) & \Longleftrightarrow\left( \exists x\right) \left(
\lnot\left( \exists y\right) \left( f(x)=y\right) \right) \\
& \Longleftrightarrow\left( \exists x\right) \left( \left( \forall y\right)
\left( f(x)\neq y\right) \right) \text{,}
\end{align*}
i, reformulant aquesta última expressió en català es té:
\enquote{Existeix un nombre real $x$ tal que per a tot nombre real $y$ es compleix $f(x)\neq y$}.

\bigskip

Finalment, passem a les regles d'inferència amb quantificadors. Hi ha
quatre regles d'inferència d'aquest tipus i, tot i que el seu ús
requereix una mica més de cura que les regles d'inferència de la
secció \ref{sec:lleis-inferencia}, s'utilitzen amb el mateix propòsit, que és
mostrar la validesa dels raonaments lògics.

\begin{description}
\item[EQU:]
\begin{tabular}{l}
$\left( \forall x\in U\right) P(x)$ \\ \hline
$P(a)$\end{tabular}
(Eliminació Quantificador Universal)

on $a$ és qualsevol element de $U$.

\item[IQU:]
\begin{tabular}{l}
$P(b)$ \\ \hline
$\left( \forall x\in U\right) P(x)$\end{tabular}
(Introducció Quantificador Universal)

on $b$ és un element arbitrari de $U$.

\item[EQE:]
\begin{tabular}{l}
$\left( \exists x\in U\right) P(x)$ \\ \hline
$P(c)$\end{tabular}
(Eliminació Quantificador Existencial)

on $c$ és un element de $U$ però el símbol \enquote{$c$} no ha d'haver
aparegut en l'argument abans.

\item[IQE:]
\begin{tabular}{l}
$P(d)$ \\ \hline
$\left( \exists x\in U\right) P(x)$\end{tabular}
(Introducció Quantificador Existencial)

on $d$ és un element de $U$.
\end{description}

En la regla EQE volem destacar el fet que $c$ no fa referència a cap símbol que ja s'ha utilitzat en l'argument. Per tant, hem de triar una
lletra nova, en lloc d'una que ja s'utilitzi per a una altra cosa.

\bigskip

Un exemple de raonament lògic senzill que implica quantificadors és
el següent:

\begin{quote}
\enquote{A cada gat simpàtic i intel·ligent li agrada el
fetge picat. Tots els gats siamesos són agradables. Hi ha un gat siamès al qual no li agrada el fetge picat. Per tant, hi ha un gat estúpid.}
\end{quote}

Formalitzem els enunciats del raonament. Per facilitar l'escriptura,
considerem que la variable $x$ té per domini el conjunt de tots els
gats. Denotem per $S$, $I$, $F$ i $G$ els predicats \enquote{és simpàtic o
agradable}, $I$, \enquote{és intel·ligent}, $F$, \enquote{agrada el
fetge picat}, i $G$, \enquote{és gat siamès}, respectivament. Aleshores, el
raonament podem simbolitzar-lo d'aquesta manera:\begin{equation*}
\begin{tabular}{ll}
$P_{1}:$ & $\left( \forall x\right) \left( S(x)\wedge I(x)\longrightarrow
F(x)\right) $ \\
$P_{2}:$ & $\left( \forall x\right) \left( G(x)\longrightarrow S(x)\right) $
\\
$P_{3}:$ & $\left( \exists x\right) \left( G(x)\wedge\lnot F(x)\right) $ \\
\hline
$C:$ & $\left( \exists x\right) \lnot I(x)$\end{tabular}
\end{equation*}

Fem ara la deducció per provar la seva validesa:

\begin{equation*}
\begin{tabular}{lll}
1. & $\left( \forall x\right) \left( S(x)\wedge I(x)\longrightarrow
F(x)\right) $ & $P_{1}$ \\
2. & $\left( \forall x\right) \left( G(x)\longrightarrow S(x)\right) $ & $P_{2}$ \\
3. & $\left( \exists x\right) \left( G(x)\wedge\lnot F(x)\right) $ & $P_{3}$
\\
4. & $G(a)\wedge\lnot F(a)$ & EQE 3 \\
5. & $G(a)\longrightarrow S(a)$ & EQU 2 \\
6. & $S(a)\wedge I(a)\longrightarrow F(a)$ & EQU 1 \\
7. & $G(a)$ & EC 4 \\
8. & $\lnot F(a)$ & EC 4 \\
9. & $S(a)$ & MP (5,7) \\
10. & $\lnot\left( S(a)\wedge I(a)\right) $ & MT (6,8) \\
11. & $\lnot S(a)\vee\lnot I(a)$ & DM 10 \\
12. & $\lnot\lnot S(a)$ & DN 9 \\
13. & $\lnot I(a)$ & ED (11,12) \\ \hline
14. & $\left( \exists x\right) \lnot I(x)$ & IQE 13\end{tabular}
\end{equation*}

Tinguem en compte que a la línia (4) hem escollit alguna lletra com \enquote{$a$}\ que no s'utilitzava abans d'aquesta línia, perquè apliquem la
regla EQE.

\bigskip

\begin{exemple}
Considereu el raonament formalitzat següent:\begin{equation*}
\begin{tabular}{ll}
$P_{1}:$ & $\left( \exists x\in U\right) \left( P(x)\wedge Q(x)\right) $ \\
$P_{2}:$ & $\left( \exists x\in U\right) M(x)$ \\ \hline
$C:$ & $\left( \exists x\in U\right) \left( M(x)\wedge Q(x)\right) $\end{tabular}
\
\end{equation*}
Es proposa la prova següent per a provar la seva validessa, és
correcte?\begin{equation*}
\begin{tabular}{lll}
1. & $\left( \exists x\in U\right) \left( P(x)\wedge Q(x)\right) $ & $P_{1}$
\\
2. & $\left( \exists x\in U\right) M(x)$ & $P_{2}$ \\
3. & $P(a)\wedge Q(a)$ & EQE 1 \\
4. & $Q(a)$ & EC 3 \\
5. & $M(a)$ & EQE 2 \\
6. & $Q(a)\wedge M(a)$ & IC (4,5) \\
7. & $\left( \exists x\in U\right) \left( M(x)\wedge Q(x)\right) $ & IQE 6\end{tabular}
\end{equation*}
$\ $
\end{exemple}

\begin{solucio}
No és correcte, perquè en el pas 5 el símbol \enquote{$a$} ha aparegut
abans.
\end{solucio}

\section{Definicions i la identitat}

Quan, per exemple, en el conjunt dels nombres enters diem que $x$ divideix
$y$, i ho escrivim $x\mid y$, volem dir que existeix un nombre enter $k$ tal
que $y=kx$. Simbòlicament:
\begin{equation*}
x\mid y\Longleftrightarrow\left(\exists k\in\mathbb{Z}\right)\left(y=kx\right).
\end{equation*}
Cal no confondre les dues formes de parlar: \enquote{$x$ divideix $y$} vol dir
$x\mid y$; en canvi, \enquote{$x$ és divisible per $y$} vol dir $y\mid x$. En tots
dos casos establim una definició del símbol \enquote{$\mid$} mitjançant termes ja
coneguts, com \enquote{nombre enter} i \enquote{producte}. De la mateixa manera, en el
conjunt dels nombres reals diem que $x\leq y$ si i només si no és el cas que
$x>y$, formalment,
\begin{equation*}
\left( \forall x,y\right) \left( x\leq y\Longleftrightarrow\lnot\left(
x>y\right) \right) \text{,}
\end{equation*}
establim la definició del símbol \enquote{$\leq$} donant el seu significat amb ajuda
del predicat ja conegut \enquote{$>$}. En aquest últim cas, per exemple, podem
substituir en qualsevol proposició el predicat \enquote{$x\leq y$} pel predicat \enquote{no
és el cas que $x>y$} sense que canviï el seu significat.

\bigskip

No volem donar una definició precisa de com cal construir correctament
una definició, però tota definició pot adoptar la forma d'una
equivalència lògica; el membre de l'esquerra ha de contenir allò
que volem definir i, el de la dreta, allò que està ja ben definit o
que el seu significat sigui comprensible immediatament, però mai hi ha
d'aparèixer el que volem definir.

\bigskip

La noció d'identitat present en molts enunciats com \enquote{$x$ és idèntic a
$y$}, \enquote{$x$ és el mateix que $y$}, o senzillament \enquote{$x$ és igual a $y$}, i que
abreugem simbòlicament per $x=y$, es pot entendre mitjançant la llei de
Leibniz: dos objectes són iguals si tenen les mateixes propietats.
Formalment, aquesta idea s'expressa així:
\begin{equation*}
x=y\Longleftrightarrow\left( \forall P\right) \left( P(x)\longleftrightarrow
P(y)\right).
\end{equation*}
Ara bé, aquesta fórmula quantifica sobre propietats o predicats $P$, no sobre
objectes del domini. Per tant, no és una fórmula de lògica de primer ordre,
sinó una formulació de segon ordre. En la lògica de primer ordre amb igualtat,
el símbol \enquote{$=$} es pren normalment com un símbol lògic bàsic, regit per les
propietats d'identitat i per la regla de substitució: si $x=y$, qualsevol
propietat que val per $x$ val també per $y$.

A matemàtiques, aquesta idea apareix constantment. Per exemple, en àlgebra
diem que dos polinomis $A(x)$ i $B(x)$ són iguals si tenen el mateix grau i
els coeficients dels monomis del mateix grau són iguals. En aquest cas,
igualar dos polinomis vol dir que podem substituir-ne l'un per l'altre en
qualsevol càlcul sense alterar-ne el resultat.

De la noció d'igualtat s'obté de manera bastant evident que
la relació d'igualtat compleix les tres propietats següents:

\begin{description}
\item[Reflexiva:] Tot objecte és igual a si mateix: $x=x$.

\item[Simètrica:] Si $x=y$, llavors $y=x$.

\item[Transitiva:] Si $x=y$ i $y=z$, llavors $x=z$.
\end{description}

Afegint al llenguatge de la lògica proposicional variables per referir-nos a
objectes, quantificadors sobre aquests objectes, predicats i, sovint, el
símbol d'igualtat, obtenim el llenguatge de la \textbf{lògica de predicats de
primer ordre}. Si, a més, permetem quantificar sobre predicats o propietats,
entrem en la \textbf{lògica de segon ordre}. La lògica de primer ordre té un
teorema de completesa molt important, però, en general, no és decidible: no hi
ha cap procediment finit que decideixi tots els raonaments possibles del seu
llenguatge.

\section{Demostracions}

Entrem a la secció més important d'aquest document que són les
demostracions a matemàtiques. És important ser conscients de les
raons per les quals hem fet fins ara una petita introducció a la lògica. N'hi ha tres raons de molt significatives:

\begin{enumerate}
\item En primer lloc, les taules de veritat que hem estudiat ens indiquen
els significats exactes de les paraules \enquote{i}, \enquote{o}, \enquote{no}, \enquote{si ... llavors ...}
i \enquote{... si i només si ...}. Per exemple, sempre que fem servir o llegim
la construcció \enquote{Si ..., llavors...} en un context matemàtic, la lògica ens indica exactament què es vol dir.

\item En segon lloc, les regles d'inferència proporcionen un camí
pel qual podem construir nova informació (proposicions) a partir
d'informació coneguda.

\item Finalment, les equivalències lògiques ens ajuden a canviar
correctament certes proposicions en unes altres proposicions amb el mateix
significat però molt més útils.
\end{enumerate}

En resum, la lògica ens ajuda a entendre els significats dels enunciats
i també a construir de nous. De fet és el llenguatge bàsic que
ens permet escriure i comprendre bé els enunciats matemàtics. Però, malgrat el seu paper fonamental, el seu lloc queda en un segon pla.
Els símbols $\wedge,\vee,\lnot,\longrightarrow,\longleftrightarrow,\forall$ i $\exists$ poques vegades s'escriuen a la pissarra o en els
llibres de text. Tot i això, hem de ser conscients dels seus significats
constantment; quan llegim o escrivim una frase relacionada amb les matemàtiques, hem d'analitzar-la amb aquests símbols, sigui mentalment o
sobre paper en brut, per tal d'entendre o comunicar bé el seu
significat, que ha de ser sempre vertader i inequívoc.

\bigskip

A matemàtiques un raonament lògic és senzillament l'enunciat
d'un teorema. La justificació que fa que un teorema sigui vàlid és la seva prova o demostració. Quan passem a la construcció de
demostracions matemàtiques, ens centrem en el contingut matemàtic de
les proposicions implicades en la prova i no ens referim explícitament
a les regles d'inferència lògica discutides en les seccions
anteriors; fer-ho seria una distracció de les qüestions matemàtiques. Tampoc utilitzem la notació lògica de les seccions
anteriors. Tot i això, farem servir les regles d'inferència de
manera implícita tot el temps, no oblidem que és el marc sobre el
qual es basa tot.

Sovint passa que la nostra intuïció ens indica el que és
important, el què creiem que pot ser cert, el què hem de provar després i més coses. Desafortunadament, els objectes matemàtics solen
ser tan complicats o abstractes que la nostra intuïció falla. Per
això construïm demostracions, per verificar que una afirmació
determinada que ens sembla intuïtivament certa és realment certa,
només així podem avançar a matemàtiques. Finalment, afegir
també que ens interessa aprendre a fer demostracions perquè ens
ajuda a entendre les nocions o fets relacionats amb el resultat que es vol
demostrar.

\bigskip

Quan es té l'enunciat d'un teorema, primer s'ha de comprendre el que
significa, segon s'ha de saber escriure amb rigor (seguint el llenguatge de
la lògica) i finalment, amb l'ajut d'altres nocions, proposicions o
teoremes, s'ha d'explorar els possibles camins per portar a terme la seva
demostració. Per construir la prova, el primer que cal fer és
especificar en rigor el què se suposa que es compleix, que anomenem les
\textbf{hipòtesis, }i després el què s'intenta demostrar, que
se'n diu \textbf{tesi}. A continuació, s'ha d'escollir una estratègia per a la prova. La següent etapa és obtenir la prova, fent ús de l'estratègia escollida. Si no es pot idear una prova amb l'estratègia escollida, potser s'hauria d'intentar una altra. No hi ha una
manera fixa de trobar la prova; sempre requereix experimentar, jugar i
provar coses diferents. En els apartats següents tractarem les maneres més importants de fer demostracions.

\subsection{Demostracions directes}

Molts teoremes tenen la forma $A\Longrightarrow B$. El camí més
senzill per provar $A\Longrightarrow B$ és fer-ho directament: suposem
que $A$ és vertadera i hem d'arribar a deduïr $B$ mitjançant
una seqüència de passos que comença en $A$ i acaba en $B$.
Aquest tipus de prova es diu \textbf{prova directa} per distinguir-la
d'altres mètodes de demostració.

\begin{exemple}
Reformuleu simbòlicament cadascun dels teoremes següents en la forma
$A\longrightarrow B$.

\begin{enumerate}
\item L'àrea d'un cercle de radi $r$ és $\pi r^{2}$.

\item Donada una recta $r$ i un punt $P$ que no és de $r$, hi ha
exactament una recta $s$ que passa pel punt $P$ i és
paral·lela a $r$.

\item En tot triangle $ABC$, els costats del qual són $a,b$ i $c$,
llavors es compleix\begin{equation*}
\frac{a}{\sin A}=\frac{b}{\sin B}=\frac{c}{\sin C}\text{.}
\end{equation*}

\item $e^{x}\cdot e^{y}=e^{x+y}$.
\end{enumerate}
\end{exemple}

\begin{solucio}
(1) Si $C(r)=$\enquote{$C$ és un cercle del pla de radi $r$}\ i $A(x)=$\enquote{\ àrea de la figura $x$}, aleshores es té $\forall r\left( C(r)\longrightarrow
A(C(r))=\pi r^{2}\right).$ De fet, si $r>0$, aleshores $C(r)=\left\{ \left(
x,y\right) \in\mathbb{R}^{2}:x^{2}+y^{2}=r^{2}\right\} \,$.

(2) Si $R(x)=$\enquote{és una recta del pla}, $P(x)=$\enquote{és un punt del pla}\ i $x\parallel y=$\enquote{x és paral·lel a
y}, aleshores\begin{equation*}
\forall x,y\left( R(x)\wedge P(y)\longrightarrow\exists z\left( R(z)\wedge
y\in z\wedge z\parallel x\right) \right)
\end{equation*}

(3) Si $T(x,y,z)=$\enquote{$T$ és un triangle del pla de costats
$x,y,z$}\ i $A(x)=$\enquote{angle oposat al costat
$x$ del triangle $T$}. Aleshores\begin{equation*}
\forall x,y,z\left( T(x,y,z)\longrightarrow\frac{x}{\sin A(x)}=\frac{y}{\sin
A(y)}=\frac{z}{\sin A(z)}\right) \text{.}
\end{equation*}
De fet, $T(x,y,z)=\left\{ (x,y,z)\in\mathbb{R}^{3}:x\leq y+z\right\}.$

(4) $\forall x,y\left( x,y\in\mathbb{R}\longrightarrow e^{x}\cdot
e^{y}=e^{x+y}\right).$
\end{solucio}

\begin{exemple}
Considerem tres nombres enters qualssevol $a,b$ i $c$. Si $a$ divideix $b$ i
$b$ divideix $c$, prova que $a$ divideix $c$.
\end{exemple}

\begin{solucio}
Primer observem que a l'enunciat hi apareix un predicat "$x$ divideix $y$",
que simbolitzem per $x\mid y$. La seva definició és: $x\mid y$ si i només si existeix $k\in\mathbb{Z}$ tal que $kx=y$. Ara hem d'escollir una estratègia per a la prova. En aquest cas escollim la
prova directa: suposem que $a\mid b$ i $b\mid c$ (hipòtesis) i hem de
veure que $a\mid c$ (tesi).

Si $a\mid b$ i $b\mid c$, llavors existeixen $k_{1},k_{2}\in\mathbb{Z}$ tals
que $k_{1}a=b$ i $k_{2}b=c$. D'aquí, per substitució es té: $k_{2}\left( k_{1}a\right) =\left( k_{1}k_{2}\right) a=c$. Per tant, existeix
$k=k_{1}k_{2}\in\mathbb{Z}$ tal que $ka=c$ i, com a conseqüència, $a\mid c$.
\end{solucio}

\begin{exemple}
Provar que per a qualssevol nombres positius $x$ i $y$ es compleix que la
mitja aritmètica és més gran que la mitja geomètrica.
\end{exemple}

\begin{solucio}
Primer observem que a l'enunciat hi apareixen dos conceptes: mitja aritmètica i geomètrica de dos nombres positius. La mitja aritmètica i geomètrica de $x$ i $y$ són respectivament: $\sqrt{xy}$ i $\dfrac{x+y}{2}
$. Escollim una prova directa: com hipòtesi suposem $x,y\geq0$, tenim
que veure $\sqrt{xy}\leq\dfrac{x+y}{2}$ (tesi).

Si $x,y\geq0$ aleshores $\sqrt{x}$ i $\sqrt{y}$ existeixen. És evident
que $(\sqrt{x}-\sqrt{y})^{2}\geq0$. Desenvolupant aquesta última expressió i passant l'arrel quadrada a l'altre costat de la desigualtat, es té\begin{align*}
x-2\sqrt{xy}+y & \geq0 \\
x+y & \geq2\sqrt{xy}
\end{align*}
Finalment, dividint tots dos costats per $2$, es té el resultat que
volíem:\begin{equation*}
\frac{x+y}{2}\geq\sqrt{xy}\text{.}
\end{equation*}
\end{solucio}

\subsection{Demostracions cap enrere}

Un altre mètode de construir una demostració és treballant
mentalment des de la tesi cap a les hipòtesis. Es diu \textbf{prova
pensada cap enrere} per distingir-la de les demés. Tot i això, la
prova finalment es construeix de forma directa. Mirem-ho construïnt la
prova cap enrere de la següent proposició:

\begin{exemple}
Provar que per a qualssevol nombres reals $x,y$ i $x<y$, es té que $4xy<(x+y)^{2}$.
\end{exemple}

\begin{solucio}
L'hipòtesi és $x,y\in\mathbb{R}$ i $x<y$, i la tesi, $4xy<(x+y)^{2}$. La prova cap enrere surt de la tesi, i, per tant, podem fer es següent:\begin{align*}
4xy & <(x+y)^{2} \\
4xy & <x^{2}+2xy+y^{2} \\
0 & <x^{2}-2xy+y^{2} \\
0 & <(x-y)^{2}
\end{align*}
De la última desigualtat s'obté que n'hi ha prou amb saber que
$x-y\neq0$, és a dir, que $x\neq y$. La hipòtesi $x<y$ garanteix
precisament aquesta condició. De fet, això que hem escrit s'havia
d'haver pensat mentalment perquè la demostració real és una prova
directa: Si $x<y$, llavors $x\neq y$ i, per tant, $x-y\neq0$. D'aquí,
es té $(x-y)^{2}>0$ i, per tant, $x^{2}-2xy+y^{2}>0$. Sumant a tots dos
costats $4xy$ es té: $x^{2}+2xy+y^{2}>4xy$ i, per tant, $(x+y)^{2}>4xy$,
com volíem demostrar.
\end{solucio}

\subsection{Demostracions per contrarecíproc}

Recordem que el contrarecíproc de $A\longrightarrow B$ és $\lnot
B\longrightarrow\lnot A$, i també l'equivalència lògica $\left(
A\longrightarrow B\right) \Longleftrightarrow\left( \lnot B\longrightarrow
\lnot A\right) $, anomenada llei del contrarecíproc. Per això, si
volem provar $A\longrightarrow B$, podem fer-ho de manera equivalent, construïnt una prova directa de $\lnot B\longrightarrow\lnot A$. L'estratègia és doncs pendre com hipòtesi que $B$ és falsa i hem
d'arribar a que $A$ és també falsa (tesi). Aquest tipus de demostració es diu \textbf{prova per contrarecíproc}. Veiem-ho amb un exemple:

\begin{exemple}
\label{ex:n2-parell}Provar que per a qualsevol nombre enter $n$, si $n^{2}$ és
parell, llavors $n$ és parell.
\end{exemple}

\begin{solucio}
Volem constuir una prova per contrarecíproc. Aleshores, les hipòtesis són que $n\in\mathbb{Z}$ i $n$ no és parell. La tesi és
que $n^{2}$ no és parell. Si $n$ no és parell, $n$ és senar i,
per tant, existeix $k\in\mathbb{Z}$ i $n=2k+1$. Hem de veure que $n^{2}$ també és senar. En efecte, com $n=2k+1$, aleshores $n^{2}=\left(
2k+1\right) ^{2}=4k^{2}+4k+1=2\left( 2k^{2}+2k\right) +1$ i com $2k^{2}+2k\in\mathbb{Z}$, s'obté que $n^{2}$ és senar, com volíem demostrar.
\end{solucio}

\subsection{Demostracions reducció a l'absurd}

El principi del tercer exclòs permet afirmar que una proposició o es
vertadera o bé falsa. Aquest principi és la base del següent mètode de demostració. En concret, si volem demostrar que una
proposició $p$ és certa, aquest mètode proposar suposar el
contrari, o sigui que $\lnot p$ és certa, i llavors per regles d'inferència hem d'arribar a una contradicció com per exemple $q\wedge\lnot
q$, on $q$ és una proposició qualsevol, per poder concluir que la
suposició que $A$ és falsa no és possible i, en conseqüència, $A$ és certa.

A matemàtiques els teoremes tenen la forma $A\longrightarrow B$ i, per
tant, hem d'adaptar el raonament anterior a aquest cas. Recordem l'equivalència lògica $\lnot\left( A\longrightarrow B\right)
\Longleftrightarrow A\wedge\lnot B$. Per tant, si volem construir una
demostració de $A\longrightarrow B$, hem de suposar que $\lnot\left(
A\longrightarrow B\right) $ és certa. Això és equivalent a
suposar que $A\wedge\lnot B$ és certa, o sigui que $A$ és certa i $B$
és falsa i d'aquí hem de provar que s'arriba a una contradicció. Una vegada s'ha arribat a la contradicció, concluïm que $A\wedge\lnot B$ és falsa i, per tant, $A\longrightarrow B$ és
vertadera.

Aquest mètode de demostració es coneix com a \textbf{prova per
reducció a l'absurd}. L'hipòtesi és $A\wedge\lnot B$ i la tesi
és una contradicció o absurd.

\begin{exemple}
Volem provar que $\sqrt{2}$ és un nombre irracional.
\end{exemple}

\begin{solucio}
Recordem que un nombre real $a$ és racional si existeixen enters $m,n$, amb
$n\neq0$, tals que $a=\dfrac{m}{n}$. Per tant, dir que $\sqrt{2}$ és
irracional equival a dir
\begin{equation*}
\lnot\left(\exists m,n\in\mathbb{Z}\right)
\left(n\neq0\wedge \sqrt{2}=\frac{m}{n}\right).
\end{equation*}
Farem una prova per reducció a l'absurd. Suposem que $\sqrt{2}$ és racional.
Aleshores existeixen enters $m,n$, amb $n\neq0$, tals que
\[
\sqrt{2}=\frac{m}{n}.
\]
Podem suposar, a més, que la fracció és irreductible. Elevem al quadrat:
\[
2=\frac{m^{2}}{n^{2}}, \qquad \text{d'on} \qquad m^{2}=2n^{2}.
\]
Així, $m^{2}$ és parell. Per l'exemple \ref{ex:n2-parell}, $m$ també és
parell, i per tant $m=2k$ per algun $k\in\mathbb{Z}$. Substituint-ho en la
igualtat anterior,
\[
4k^{2}=2n^{2}, \qquad \text{i per tant} \qquad n^{2}=2k^{2}.
\]
D'aquí deduïm que $n^{2}$ és parell i, novament, que $n$ és parell. Això
contradiu que la fracció $\dfrac{m}{n}$ fos irreductible, perquè $m$ i $n$
tindrien el factor comú $2$. Per tant, la suposició inicial és falsa i
$\sqrt{2}$ és irracional.
\end{solucio}

\begin{exemple}
Volem provar que no existeix cap nombre real $x$ tal que $x^{2}+1=0$.
\end{exemple}

\begin{solucio}
Primer, escrivim simbòlicament l'enunciat: $\lnot\left( \exists x\in\mathbb{R}\right) \left( x^{2}+1=0\right).$ Veiem que és la negació
d'una proposició existencial. Fent ús de la següent equivalència lògica\begin{equation*}
\lnot\left( \exists x\in\mathbb{R}\right) \left( x^{2}+1=0\right)
\Longleftrightarrow\left( \forall x\in\mathbb{R}\right) \left(
x^{2}+1\neq0\right) \text{,}
\end{equation*}
i de la regla d'inferència IQU:\begin{equation*}
\begin{tabular}{l}
$P(b)$ \\
$\left( \forall x\in U\right) P(x)$\end{tabular}
\ \text{,}
\end{equation*}
llavors només cal provar que prenent un nombre real qualsevol $b$, es
compleix que $b^{2}+1\neq0$. Però, això és immediat perquè $b^{2}\geq0$ i, per tant, $b^{2}+1\geq1$; en particular, $b^{2}+1\neq0$.
\end{solucio}

\subsection{Demostracions per contraexemple}

A matemàtiques també trobem enunciats amb quantificadors que
rebutjant propietats. Però, de fet rebutjar una propietat enunciada per
una proposició \enquote{$p$} és el mateix que provar \enquote{$\lnot p$}. Per
exemple, per demostrar que no és el cas que $\left( \forall x\in
U\right) P(x)$, primer cal pensar en la següent equivalència lògica:\begin{equation*}
\lnot\left( \forall x\in U,P(x)\right) \Longleftrightarrow\exists x\in
U,\lnot P\left( x\right)
\end{equation*}
i després recordar la regla d'inferencia IQE:\begin{equation*}
\begin{tabular}{l}
$\lnot P(d)$ \\
$\left( \exists x\in U\right) \lnot P(x)$\end{tabular}
\text{.}
\end{equation*}
Llavors, només cal trobar un exemple pel qual $P(x)$ és fals.
Aquesta manera de fer la demostració és coneix com \textbf{prova per
contraexemple}.

\begin{exemple}
Volem provar que no és veritat que la suma de dos nombres irracionals
és irracional.
\end{exemple}

\begin{solucio}
Per buscar una estratègia, primer haurem de reformular l'enunciat lògicament: si $\mathbb{I}=\mathbb{R}\smallsetminus\mathbb{Q}$ és el
conjunt dels nombres irracionals, aleshores\begin{equation*}
\lnot\left( \forall x,y\in\mathbb{I}\right) \left( x+y\in\mathbb{I}\right)
\end{equation*}
Hem vist que aquest enunciat és equivalent al següent:\begin{equation*}
\left( \exists x,y\in\mathbb{I}\right) \left( x+y\notin\mathbb{I}\right)
\end{equation*}
Per tant, per demostrar l'enunciat només cal trobar dos nombres
irracionals la suma dels quals no és irracional. Per exemple, $\sqrt{2}$
i $-\sqrt{2}$ són irracionals, però
\begin{equation*}
\sqrt{2}+(-\sqrt{2})=0,
\end{equation*}
i $0$ és racional. Aquest és el contraexemple buscat.
\end{solucio}

\subsection{Demostracions per casos}

A vegades trobem enunciats que són de la forma $\left( C\vee D\right)
\longrightarrow B$, és a dir, que inclouen una disjunció en
l'antecedent. En aquests casos, la següent equivalència
\begin{equation*}
\left( C\vee D\right) \longrightarrow B\Longleftrightarrow\left(
C\longrightarrow B\right) \wedge\left( D\longrightarrow B\right) \text{,}
\end{equation*}
proporciona l'estratègia que hem de seguir. En efecte, per demostrar que
$\left( C\vee D\right) \longrightarrow B$ és certa és suficient
demostrar que $C\longrightarrow B$ i $D\longrightarrow B$ són ambdues
certes. De fet, la demostració la construïm per casos (en aquest
cas, n'hi ha dos casos) segons la forma que tingui l'antecedent. La
demostració de cadascun dels casos seguirà l'estratègia que faci
falta. Aquest mètode de demostració es coneix com \textbf{prova per
casos}.

\begin{exemple}
Volem provar que si $n$ és un nombre enter, llavors $n^{2}+n$ és
parell.
\end{exemple}

\begin{solucio}
Sabem que tot nombre enter $n$ és o parell o bé senar. Aquesta idea
proporciona l'estratègia de la demostració: construïr una prova
per casos. Volem provar que (1) si $n$ és parell, aleshores $n^{2}+n$
és parell, i (2) si $n$ és senar, aleshores $n^{2}+n$ és parell.

\begin{description}
\item[Cas 1:] Si $n$ és parell, llavors per definició existeix $k\in\mathbb{Z}$ tal que $n=2k$. Aleshores es té:
\begin{equation*}
n^{2}+n=4k^{2}+2k=2\left( 2k^{2}+k\right) \text{,}
\end{equation*}
i, per tant, $n^{2}+n$ és parell perquè $2k^{2}+k\in\mathbb{Z}$.

\item[Cas 2:] Si $n$ és senar, llavors per definició existeix $m\in\mathbb{Z}$ tal que $n=2m+1$. Aleshores es té:\begin{align*}
n^{2}+n & =\left( 2m+1\right) ^{2}+2m+1 \\
& =4m^{2}+6m+2=2\left( 2m^{2}+3m+1\right) \text{,}
\end{align*}
i, per tant, $n^{2}+n$ és parell perquè $2m^{2}+3m+1\in\mathbb{Z}$.
\end{description}
\end{solucio}

\bigskip

Ara passem als teoremes de la forma $A\longrightarrow\left( C\vee D\right) $
on hi trobem una disjuntiva en el conseqüent. Una estratègia per a
fer la demostració és fent ús primer de la llei del contrarecíproc i després una de les lleis de De Morgan:
\begin{align*}
A & \longrightarrow\left( C\vee D\right) \Longleftrightarrow\lnot\left(
C\vee D\right) \longrightarrow\lnot A \\
& \Longleftrightarrow\left( \lnot C\wedge\lnot D\right) \longrightarrow
\lnot A
\end{align*}
A partir d'aquí, construïm la prova directa o per reducció al
absurd. Veiem-ho en un exemple.

\begin{exemple}
Considerem dos nombres reals $x$ i $y$. Volem provar que si $xy$ és
irracional, llavors $x$ o $y$ és irracional.
\end{exemple}

\begin{solucio}
Segons el que hem vist anteriorment, aquest enunciat és equivalent al següent: Si $x$ és racional i $y$ també ho és, llavors $xy$
és racional. Però això és evident perquè sabem que la
multiplicació és una operación interna del cos dels racionals.
\end{solucio}

\bigskip

Després d'haver discutit sobre l'aparició de la connectiva $\vee$,
anem ara a tractar els casos quan aprareix la connectiva $\wedge$ en els
teoremes de la forma $A\longrightarrow B$. Per demostar un enunciat de la
forma $A\longrightarrow\left( C\wedge D\right) $ fem ús de l'equivalència següent:\begin{equation*}
A\longrightarrow\left( C\wedge D\right) \Longleftrightarrow\left(
A\longrightarrow C\right) \wedge\left( A\longrightarrow D\right)
\end{equation*}
A partir d'aquí, la demostració consisteix en provar que $A\longrightarrow C$ i $A\longrightarrow D$ són certes.

Si l'enunciat és de la forma $\left( C\wedge D\right) \longrightarrow B$, llavors l'estratègia és prendre com hipòtesis que $C$ i $D$ són certes i com a tesi que $B$ és certa. La prova pot ser directe o
per reducció a l'aburd.

Finalment, volem tractar una altra forma lògica que presenten molts
teoremes. Són els teoremes de la forma $A\longleftrightarrow B$ i que
ens donen condicions necessàries i suficients. La llei del bicondicional
ens dona la resposta per construir proves d'aquest teoremes:\begin{equation*}
A\longleftrightarrow B\Longleftrightarrow\left( A\longrightarrow B\right)
\wedge\left( B\longrightarrow A\right) \text{.}
\end{equation*}
Per tant, és suficient provar $A\longrightarrow B$ i $B\longrightarrow A$, cadascun dels quals es pot demostrar mitjançant qualsevol dels mètodes que hem vist fins ara.

\begin{exemple}
Volem provar que si $x$ i $y$ són dos nombres reals, llavors $xy=0$ si i
només si $x=0$ o $y=0$.
\end{exemple}

\begin{solucio}
L'enunciat expressat formalment és:
\begin{equation*}
\left( \forall x,y\in\mathbb{R}\right) \left( xy=0\longleftrightarrow
(x=0\vee y=0)\right).
\end{equation*}
Fixem dos nombres reals arbitraris $x$ i $y$ i provem les dues implicacions.

Primer, suposem que $xy=0$. Volem provar que $x=0$ o $y=0$. Ho fem pel
contrarecíproc: si $x\neq0$ i $y\neq0$, aleshores existeixen els inversos
$x^{-1}$ i $y^{-1}$. Si fos $xy=0$, multiplicant pels inversos obtindríem
\[
(x^{-1}y^{-1})(xy)=x^{-1}y^{-1}\cdot0,
\]
és a dir, $1=0$, una contradicció. Per tant, si $xy=0$, necessàriament
$x=0$ o $y=0$.

Recíprocament, si $x=0$ o $y=0$, aleshores $xy=0$ de manera immediata: si
$x=0$, $xy=0\cdot y=0$; i si $y=0$, $xy=x\cdot0=0$.

Per tant, $xy=0$ si i només si $x=0$ o $y=0$.
\end{solucio}

\subsection{Demostracions d'existència i unicitat}

És molt freqüent trobar demostracions que impliquin l'existència
d'un objecte que compleix una determinada propietat i també la seva
unicitat. Quant a la prova d'existència només cal trobar l'objecte
en concret que compleix la propietat. Quant a la unicitat, l'estratègia més comuna és suposar que $x$ i $y$ són dos objectes que satisfan
la propietat donada, llavors hem de demostrar que $x=y$.

\begin{exemple}
Volem provar que la proporció entre dos nombres reals $a$ i $b$, amb
$a>b>0$, que compleixen
\begin{equation*}
\frac{a}{b}=\frac{a+b}{a}
\end{equation*}
és un nombre irracional i, a més, és únic.
\end{exemple}

\begin{solucio}
Posem
\[
x=\frac{a}{b}.
\]
Com que $a>b>0$, tenim $x>1$. La condició de l'enunciat es transforma en
\begin{equation*}
x=\frac{x+1}{x},
\end{equation*}
i, per tant,
\begin{equation}
x^{2}-x-1=0.
\label{eq:phi-quadratica}
\end{equation}
Resolem l'equació de segon grau:
\[
x=\frac{1\pm\sqrt{5}}{2}.
\]
Com que $x>1$, l'única solució admissible és
\begin{equation*}
\Phi=\frac{1+\sqrt{5}}{2}.
\end{equation*}
Això prova la unicitat de la proporció: qualsevol parella $a>b>0$ que
satisfaci la propietat ha de tenir el mateix quocient $a/b=\Phi$.

També existeix una parella que dona aquesta proporció: per exemple, podem
prendre $b=1$ i $a=\Phi$. Aleshores $a>b>0$ i, com que $\Phi$ satisfà
\eqref{eq:phi-quadratica}, es compleix
\[
\Phi=\frac{\Phi+1}{\Phi},
\]
que és exactament la propietat demanada.

Finalment, provem que $\Phi$ és irracional. Ja sabem, pel mateix argument
utilitzat per demostrar que $\sqrt{2}$ és irracional, que $\sqrt{5}$ és
irracional. Si $\Phi$ fos racional, aleshores
\[
\sqrt{5}=2\Phi-1
\]
també seria racional, perquè els racionals són tancats per suma, resta i
multiplicació per enters. Això és una contradicció. Per tant, $\Phi$ és
irracional.
\end{solucio}

\subsection{Demostracions per inducció\label{subsec:induccio}}

La inducció matemàtica és una tècnica útil per demostrar
enunciats sobre nombres naturals. Considerem, per exemple, l'enunciat següent: "Si $n$ és parell, llavors $n^{2}$ és divisible per 4".
Volem provar que aquest enunciat és correcte. Com ho fem? Aplicant
inducció sobre $n$, que vol dir provar els dos passos següents:

\begin{enumerate}
\item Primer hem de provar que la propietat és vertadera pel primer
element. En aquest cas el primer element és $2$ perquè tractem amb
nombres naturals parells. La propietat es certa perquè $2^{2}=4$ que
és divisible per 4.

\item Per a tot nombre natural $n$ (no sabem quin és), si la propietat
és certa per $n$, llavors hem de provar que també ho és per el
següent element: Suposem \ $n$ és parell i que $n^{2}$ és
divisible per 4, aleshores hem de provar que el següent parell $(n+2)^{2} $ és divisible per 4. En efecte, $(n+2)^{2}=n^{2}+4n+4=4k+4n+4=4(k+n+1)$, on $n^{2}=4k$ i $k\in\mathbb{N}$.
\end{enumerate}

Concluïm que la propietat és vàlida per a tot nombre natural
parell.

Intuïtivament, aquestes dues condicions permeten assegurar que la
propietat es compleix per tots els naturals. En efecte, $P(2)$ és
vertadera, $P(2)\longrightarrow P(4)$ és vertadera i, per tant, $P(4)$
també ho és;\ $P(4)\longrightarrow P(6)$ és certa i, per tant, $P(6)$ també. I així succesivament.

\bigskip

La base d'aquest mètode de demostració és el principi d'inducció que, de fet és un dels axiomes que defineixen els nombres
naturals. Aquest principi diu: Si $P(n)$ és un enunciat sobre els
naturals i es compleixen les dues condicions següents: A la primera
part, anomenada \textbf{cas base}, mostrem que $P(1)$ es compleix (el
primer element no necessàriament és 1, depèn de cada cas). A la
segona part, anomenada pas inductiu, suposem que $n$ és un natural tal
que $P(n)$ és certa, tot i que no sabem què és $n$, i hem de deduïr que $P(n+1)$ és certa. La suposició que $P(n)$ és
vertadera es diu \textbf{hipòtesi d'inducció}. La conclusió
d'aquest raonament és que $P(n)$ és vàlida per a tot $n$
natural. Aquest mètode de demostració es coneix com \textbf{prova
per inducció} sobre $n$.

\bigskip

Veiem en detall aquest mètode en un exemple:

\begin{exemple}
Volem demostrar que per a tot $n$ natural es compleix\begin{equation*}
1+2+\cdots+n=\frac{n(n+1)}{2}\text{.}
\end{equation*}
\end{exemple}

\begin{solucio}
Construïm una prova d'inducció sobre $n$, sent $P(n)=$ \enquote{$1+2+\cdots
+n=\dfrac{n(n+1)}{2}$}.

\begin{description}
\item[Cas base:] $P(1)$ és $1=\frac{1(1+1)}{2}$ i, evidentment és
veritat.

\item[Hipòtesi d'inducció:] Suposem que $P(n)$ és certa, o sigui
que $1+2+\cdots+n=\dfrac{n(n+1)}{2}$.

\item[Tesi:] Hem de demostrar que $P(n+1)=$\enquote{$1+2+\cdots+n+(n+1)=\dfrac{(n+1)(n+2)}{2}$} és certa.
\end{description}

En efecte, aplicant l'hipòtesi d'inducció, es té\begin{align*}
1+2+\cdots+n+(n+1) & =\dfrac{n(n+1)}{2}+n+1 \\
& =\frac{\left( n+1\right) (n+2)}{2}
\end{align*}
que és el que volíem veure.

\begin{description}
\item[Conclusió:] Pel principi d'inducció sobre $n$, deduïm que
$P(n)$ és certa per a tot $n\in\mathbb{N}$.
\end{description}
\end{solucio}

\begin{exemple}
Volem demostrar que per a tot nombre natural $n$ es compleix que $8^{n}-3^{n} $ és divisible per 5.
\end{exemple}

\begin{solucio}
El cas base és per $n=1$: $8-3=5$ que és divisible per $5$. Suposem
ara que, per a un cert nombre natural $n$, $8^{n}-3^{n}$ és
divisible per 5 (hipòtesi d'inducció). Hem de demostrar que $8^{n+1}-3^{n+1}$ també és divisible per 5. En efecte, tenim\begin{align*}
8^{n+1}-3^{n+1} & =8\cdot8^{n}-3\cdot3^{n} \\
& =8\cdot8^{n}-\left( 8-5\right) \cdot3^{n} \\
& =8\cdot\left( 8^{n}-3^{n}\right) +5\cdot3^{n}
\end{align*}
però, per hipòtesi d'inducció, $8^{n}-3^{n}=5k$, sent $k\in\mathbb{N}$. Per tant,\begin{equation*}
8^{n+1}-3^{n+1}=5\cdot\left( 8k+3^{n}\right)
\end{equation*}
i, d'aquí, s'obté que $8^{n+1}-3^{n+1}$ és divisible per 5 perquè $8k+3^{n}\in\mathbb{N}$. Com a conseqüència, $8^{n}-3^{n}$
és divisible per 5 per a tot nombre natural $n$.
\end{solucio}

Les definicions \textbf{inductives} estan implícites en les definicions
de diverses funcions molt comunes que impliquen nombres enters no negatius.
Per exemple, el factorial d'un nombre enter no negatiu $n$, designat per $n!$, es defineix inductivament d'aquesta manera:

\begin{itemize}
\item $0!=1;$

\item $\left( n+1\right) !=\left( n+1\right) \cdot n!$.
\end{itemize}

Un altre exemple són les definicions per \textbf{recursió}. La successió de Fibonacci és un bon exemple: $1,1,2,3,5,8,13,...$, que podem
definir per recursió d'aquesta manera:

\begin{itemize}
\item $u_{1}=1$,

\item $u_{2}=1$,

\item $u_{k+1}=u_{k-1}+u_{k}$ per $k\in\mathbb{N}$ i $k\geq2$.
\end{itemize}

\section{Teories axiomàtiques}

Resulta que, per demostrar alguna cosa, es necessita el coneixement
d'algunes veritats anteriors. La lògica només proporciona les
maneres que podem deduir una afirmació d'altres, però necessitem
algunes afirmacions per començar. Aquestes afirmacions inicials
s'anomenen \textbf{axiomes} o \textbf{postulats} d'ençà que Euclides
va escriure els seus Elements cap al 300 aC. Per exemple, el primer postulat
d'Euclides diu \enquote{donats dos punts, existeix una única línea recta
que passa per ells}, enunciat que la majoria de la gent considera evident.
Tot i això, en l'enfocament axiomàtic modern els axiomes no es veuen
necessàriament com a veritats evidents per si mateixos, sinó
simplement com a afirmacions que suposem que són certes. Fem matemàtiques explorant el que es deriva de la veritat d'aquests axiomes mitjançant les regles de deducció de la lògica.

\bigskip

Les matemàtiques modernes es basen en la teoria de conjunts i la lògica. La majoria dels objectes matemàtics, com ara punts, línies,
nombres, funcions, successions, grups, etc., són realment conjunts. Per
tant, és necessari començar a coneixer els axiomes de la teoria de
conjunts. Aquí no tractarem aquest punt però es pot trobar en
qualsevol text de teoria de conjunts de nivell avançat.

A més, tenim els axiomes que defineixen cadascuna de les teories. Per
exemple, com hem comentat abans, en geometria euclidiana tenim l'axioma que
afirma que existeix una i només una recta que passa per dos punts
diferents. En les matemàtiques modernes no té cap sentit discutir
aquest axioma. Si no s'assumeix, simplement no podem anomenar al objectes
amb els noms \enquote{línia recta} i \enquote{punt}. Els axiomes serveixen aquí
com a definicions que caracteritzen aquest sistema matemàtic. De fet,
una teoria matemàtica és com un joc d'escacs i els axiomes
corresponen a les regles del joc. Si no accepteu una regla, el joc deixa de
ser escacs i és un altre cosa.

\bigskip

Les matemàtiques s'ocupen de sistemes abstractes de diversos tipus,
definits cadascun per un conjunt adequat d'axiomes, que serveixen per
caracteritzar la seva estructura. Però tot i que, des del punt de vista
de les matemàtiques pures, cada estructura es considera autònoma,
l'esquema matemàtic sol tenir una o més realitzacions concretes;
és a dir, l'estructura sol trobar-se (possiblement només fins a un
cert grau d'aproximació) en un sistema més concret. La geometria
euclidiana abstracta de tres dimensions, per exemple, té una de les
seves realitzacions l'estructura de l'espai ordinari.

\bigskip

Una teoria està formada per tots els teoremes que es dedueixen dels seus
axiomes. Com a exemples, presentem dues teories axiomàtiques clàssiques de les matemàtiques, i completarem aquest document introductori
fent de cadascuna d'elles la demostració d'un teorema.

\subsection{Teoria axiomàtica de la geometria plana}

La geometria euclidiana plana es pot considerar com el conjunt de resultats
que s'obtenen a partir dels cinc postulats d'Euclides per deducció lògica. Com es pot observar de seguida, els postulats d'Euclides estan pensats
per fer geometria en el pla fent ús de la regla i el compàs. Aquests
postulats són:

\begin{enumerate}
\item Es pot traçar una línia recta entre dos punts qualsevol i el
resultat és un segment de línia recta.

\item Qualsevol segment de línia recta es pot ampliar indefinidament.

\item Donat un punt i un segment de línia de recta que comencen pel
punt, podeu dibuixar un cercle centrat en el punt donat amb el segment de línia donat com a radi.

\item Tots els angles rectes són iguals.

\item Si dues rectes en un pla es troben amb una altra recta i si la suma
dels angles interns d'un costat és inferior a dos angles rectes, les
rectes es reuniran si s'estenen prou al costat on es suma la suma dels
angles és inferior a dos angles rectes.

\FRAME{dtbpF}{10.4493cm}{7.9693cm}{0pt}{}{}{logi1.jpg}{\special{ language "Scientific Word"; type "GRAPHIC";
maintain-aspect-ratio TRUE; display "USEDEF"; valid_file "F"; width
10.4493cm; height 7.9693cm; depth 0pt; original-width 10.3768cm;
original-height 7.8969cm; cropleft "0"; croptop "1"; cropright "1";
cropbottom "0"; filename 'img/logi1.jpg';file-properties "XNPEU";}}
\end{enumerate}

El cinquè axioma és conegut com l'\textbf{axioma de les
paral·leles} quan és reformulat d'aquesta manera:
\enquote{Donada qualsevol recta i un punt que no estigui sobre aquesta recta,
existeix una única recta que passa pel punt i és
paral·lela a la recta donada}. De fet, la geometria clàssica euclidiana es distingeix per aquest axioma. Si no acceptem aquest
resultat, surt per exemple l'\textbf{axioma de Lobachevsky}: Donada una
recta i un punt exterior a ella, existeixen almenys dues rectes que passen
per aquest punt i que no tallen la recta donada. La geometria caracteritzada
pels quatre primers postulats d'Euclides i l'axioma de Lobachevsky
distingeix la geometria anomenada hiperbòlica.

\begin{teorema}[Teorema de Pitàgores]
En tots els triangles rectangles es compleix que el quadrat del costat
oposat a l'angle recte és igual a la suma dels quadrats dels costats que
comprenen a l'angle recte.
\end{teorema}

\begin{prova}
En efecte, considerem un triangle rectangle qualsevol $ABC$, on hi suposem
l'angle $\angle BAC$ recte. Podem construir un quadrat sobre cada segment
del triangle, i obtenir d'aquesta manera una construcció com la de la
figura següent. \FRAME{dtbpF}{8.4636cm}{9.7969cm}{0pt}{}{}{logi2.jpg}{\special{ language "Scientific Word"; type "GRAPHIC"; maintain-aspect-ratio
TRUE; display "USEDEF"; valid_file "F"; width 8.4636cm; height 9.7969cm;
depth 0pt; original-width 8.3933cm; original-height 9.7267cm; cropleft "0";
croptop "1"; cropright "1"; cropbottom "0"; filename
'img/logi2.jpg';file-properties "XNPEU";}}

Observem primer que els triangles $DCB$ i $ABI$ són iguals, perquè $AB=BD$, $BI=BC$ i $\angle DBC=\angle ABI$. Llavors, l'àrea del quadrat $ABDE$ és el doble de l'àrea del triangle $DCB$, perquè les dues
figures tenen la mateixa base i es troben entre les mateixes
paral·leles. Per la mateixa raó l'àrea del
rectangle $BIJK$ és el doble de l'àrea del triangle $ABI$.
D'aquestes dues conclusions, obtenim que l'àrea del rectangle $BIJK$
és igual que l'àrea del quadrat $ABDE$. Anàlogament, s'obté
que l'àrea del rectangle $CHKJ$ és igual que l'àrea del quadrat $ACGF$. Per tant, com l'àrea del quadrat $BIHC$ és igual a la suma de
les àrees dels rectangles $BIKJ$ i $CHKJ$, per les igualtats demostrades
obtenim que l'àrea del quadrat $BIHC$ de catet oposat a l'angle recte,
és igual a la suma de les àrees dels quadrats $ABDE$ i $ACGF$, que
tenen per catets els segments que comprenen l'angle recte.
\end{prova}

\subsection{Teoria axiomàtica de l'aritmètica}

A partir dels postulats de Peano, és possible demostrar totes les
propietats esperades dels nombres naturals i, a partir dels nombres naturals,
és possible construir primer els enters, després els nombres
racionals i després els nombres reals. El conjunt dels nombres naturals
és el conjunt $\mathbb{N}$, l'existència del qual ve donada pels
postulats de Peano que són:

\begin{enumerate}
\item Existeix un element $1$ que pertany al conjunt dels nombres naturals $\mathbb{N}$.

\item Existeix una funció $s:\mathbb{N}\longrightarrow\mathbb{N}$ tal
que compleix les propietats següents:

\begin{enumerate}
\item No hi ha cap nombre natural $n$ tal que $s(n)=1.$

\item Per a tots $m,n\in\mathbb{N}$, si $s(m)=s(n)$, llavors $m=n$.

\item Per a cada subconjunt $K\subseteq\mathbb{N}$, si $1\in K$ i per a cada
nombre natural $n\in K$ es té que $s(n)\in K$, llavors tot nombre
natural és de $K$.
\end{enumerate}
\end{enumerate}

Si pensem intuïtivament en la funció $s$ com la que assigna a cada
nombre natural el seu \textbf{següent}, llavors la part (a) dels
postulats diu que $1$ és el \textbf{primer element} en $\mathbb{N}$,
perquè no és el successor de res. La part (c) es coneix com a principi
d'inducció i constitueix la base de les demostracions per inducció
que hem tractat en l'apartat \ref{subsec:induccio}. Més endavant veurem que
d'aquest principi es pot deduir que qualsevol subconjunt no buit de
$\mathbb{N}$ té primer element; aquesta propietat s'anomena bon ordre.

\bigskip

La pregunta que ens fem de seguida és com sabem que hi ha un conjunt i
un element del conjunt i una funció del conjunt en si mateix, que
satisfan els postulats de Peano? La resposta no la tractarem aquí. Com
hem assenyalat més amunt la teoria de conjunts és la base de les \
matemàtiques modernes, llavors no és necessari suposar
addicionalment que els postulats de Peano es compleixen, perquè l'existència d'alguna cosa que compleixi els postulats de Peano es pot deduir
dels axiomes de la teoria de conjunts.

\bigskip

Per a completar aquest apartat proposem demostrar una propietat característica del conjunt dels nombres naturals: tot subconjunt $K$ no buit de
$\mathbb{N}$ està ben ordenat, és a dir, que $K$ té primer
element.

\begin{teorema}
Qualsevol conjunt no buit de nombres naturals té un primer element.
\end{teorema}

\begin{prova}
La demostració la farem per reducció a l'absurd utilitzant alhora el
principi d'inducció: Considerem un conjunt no buit de nombres naturals $K$
sense cap primer element. Provarem per inducció que cap nombre natural no
pertany a $K$. El cas base és $1\notin K$: si $1\in K$, com que $1$ és
el primer element de $\mathbb{N}$, també seria el primer element de $K$,
contra la hipòtesi. Suposem ara que $1,2,\ldots,n$ no pertanyen a $K$ i
provem que $n+1\notin K$. Si $n+1\in K$, aleshores cap nombre natural més
petit que $n+1$ pertany a $K$, i per tant $n+1$ seria el primer element de
$K$, contradient novament la hipòtesi. Així, per inducció, cap
nombre natural no pertany a $K$, és a dir, $K=\emptyset$, en contradicció
amb el fet que $K$ era no buit. Per tant, $K$ té primer element.
\end{prova}